% Lesson 34 — Grammar: More difficult structures with adjectives: How + adj, the superlative, too, to, etc — Anglais Tle A — Cours signé OnBuch+
% Programme officiel MINESEC. Compiler avec : tectonic EN34-grammar-more-difficult-structures-with-adjectives-how-a.tex
\newcommand{\DOCMATIERE}{Anglais}
\newcommand{\DOCNIVEAU}{Tle A}
\newcommand{\DOCMODULE}{Chapter 3 : Environment, Well-being and Health}
\newcommand{\DOCLECON}{EN34}
\newcommand{\DOCTITRE}{Lesson 34 — Grammar: More difficult structures with adjectives: How + adj, the superlative, too, to, etc}
% OnBuch+ — préambule « cours signés » ENRICHI (compilé avec tectonic / XeLaTeX)
% Système visuel « Pop » d'OnBuch+ : fond crème (jamais blanc), bordures encre
% épaisses, ombres « dures » décalées, titres Archivo Black en capitales.
% Version MATHÉMATIQUES : graphes (pgfplots/tikz), boîtes pédagogiques
% (définition, propriété, méthode, attention, le savais-tu, exercice résolu,
% exercice, TP, bilan), page de garde et signature OnBuch+ ancrée en bas.
%
% Macros attendues dans main.tex AVANT \input{preamble} :
%   \DOCMATIERE  \DOCNIVEAU  \DOCMODULE  \DOCLECON (numéro)  \DOCTITRE
\documentclass[11pt]{article}

\usepackage{fontspec}
\usepackage[english,french]{babel} % français = langue principale ; l'anglais via \eng{..} / textebox
\usepackage[a4paper,margin=2cm,top=3.4cm,bottom=2.8cm,headheight=1.4cm,footskip=1.3cm]{geometry}
\usepackage{amsmath,amssymb,amsfonts}
\usepackage{mathtools} % \xmapsto, \coloneqq... souvent utilisés par les modèles
\usepackage{cancel} % simplifications barrées (bilans de forces)
% \abs{z} (module, valeur absolue) et \norm{u} (norme), utilisés par les modèles
\providecommand{\abs}{}\let\abs\relax\DeclarePairedDelimiter{\abs}{\lvert}{\rvert}
\providecommand{\norm}{}\let\norm\relax\DeclarePairedDelimiter{\norm}{\lVert}{\rVert}
\usepackage[table]{xcolor} % [table] : fournit aussi \rowcolors (alternance de lignes)
\usepackage{pagecolor}
\usepackage{tikz}
\usetikzlibrary{arrows.meta,positioning,calc,shapes.geometric,shapes.misc,shapes.arrows,decorations.pathmorphing,decorations.pathreplacing,decorations.markings,patterns,shadows,mindmap,trees,fit,backgrounds,matrix,intersections,angles,quotes}
\usepackage{pgfplots}
\usepackage[version=4]{mhchem} % équations nucléaires \ce{^{235}_{92}U}
\usepackage[european,siunitx]{circuitikz} % schémas électriques
\pgfplotsset{compat=1.18}
\usepgfplotslibrary{fillbetween,groupplots} % aires entre courbes (calcul intégral)
\usepackage{enumitem}
\usepackage{fancyhdr}
% [most] charge toutes les bibliothèques tcolorbox (listings, minted, raster,
% documentation...) et gonfle inutilement la mémoire TeX sur un document long
% (~300 boîtes) : on ne charge que ce qui est réellement utilisé.
\usepackage[skins,breakable]{tcolorbox}
\usepackage{graphicx}
\usepackage{colortbl}
\usepackage{booktabs}
\usepackage{tabularx}
\usepackage{diagbox} % case d'en-tête diagonale des tableaux à double entrée
% Colonnes extensibles centrée / gauche / droite (C, L, R), souvent utilisées par les modèles
\newcolumntype{C}{>{\centering\arraybackslash}X}
\newcolumntype{L}{>{\raggedright\arraybackslash}X}
\newcolumntype{R}{>{\raggedleft\arraybackslash}X}
\usepackage{array}
\usepackage{multicol}
\usepackage{multirow}
\usepackage{siunitx}
% parse-numbers=false : accepte aussi \SI{2,5 \times 10^{-3}}{...}
\sisetup{locale=FR,output-decimal-marker={,},group-minimum-digits=4,parse-numbers=false}
\DeclareSIUnit\ohm{\ensuremath{\symup{\Omega}}} % U+2126 absent de la police
\DeclareSIUnit\division{div} % réglages d'oscilloscope (ms/div, V/div)
\DeclareSIUnit\tr{tr} % tours (vitesse de rotation, ex. \SI{1500}{\tr\per\minute})
\DeclareSIUnit\molar{M} % concentration molaire (ex. \SI{3}{\molar}, \SI{1}{\micro\molar})
\DeclareSIUnit\an{an} % année (le modèle confond parfois avec \year, un compteur TeX) — ex. \SI{5}{\kilo\meter\per\an}
\DeclareSIUnit\FCFA{FCFA} % franc CFA (ex. \SI{500}{\FCFA\per\kilo\gram})
\DeclareSIUnit\degreeBrix{\textdegree Brix} % teneur en sucre d'un sirop/jus (réfractométrie)
\DeclareSIUnit\month{mois} % le modèle utilise parfois \month, un compteur TeX, comme unité de durée
\usepackage{qrcode}
% \degree : commande fréquemment utilisée par les modèles pour un angle ou une
% température en dehors de \SI{}{} (non fournie par les paquets chargés).
\providecommand{\degree}{^{\circ}}
% \qed (fin de démonstration), utilisé par les modèles sans amsthm
\providecommand{\qed}{\hfill$\square$}
% \barre : double barre des valeurs interdites dans un tableau de variations (tkz-tab)
\providecommand{\barre}{\ensuremath{\|}}
% environnement proof (amsthm non chargé) utilisé par les modèles
\ifdefined\proof\else\newenvironment{proof}[1][Démonstration]{\par\noindent\textit{#1.}\ }{\qed\par}\fi
\usepackage{lastpage}
\usepackage{hyperref}
\hypersetup{hidelinks}

% ============================================================
% Palette « Pop » OnBuch+ — mêmes hex que la classe P (plus_ui.dart)
% ============================================================
\definecolor{poporange}{HTML}{F59321}
\definecolor{popdark}{HTML}{E07A0C}
\definecolor{popcream}{HTML}{FFF6E8}
\definecolor{popink}{HTML}{1C1714}
\definecolor{poppurple}{HTML}{7A5AE0}
\definecolor{popgreen}{HTML}{2E9E6A}
\definecolor{poppink}{HTML}{E0457B}
\definecolor{popblue}{HTML}{2F7DE1}
\definecolor{popgold}{HTML}{C98A12}
\definecolor{popmuted}{HTML}{8A7F76}
\definecolor{popline}{HTML}{EADFD2}
\definecolor{popgray}{HTML}{8A7F76}
\colorlet{popgrey}{popgray}
% Alias tolérants (le modèle nomme parfois les couleurs différemment)
\colorlet{popwhite}{white}
\colorlet{popblack}{popink}
\colorlet{popred}{poppink}
\colorlet{popyellow}{popgold}
% Teintes claires (fonds de tableaux/encadrés) — le modèle nomme parfois
% les couleurs avec un suffixe L (ex. popblueL) sans les avoir définies.
\colorlet{poporangeL}{poporange!20}
\colorlet{popdarkL}{popdark!20}
\colorlet{poppurpleL}{poppurple!20}
\colorlet{popgreenL}{popgreen!20}
\colorlet{poppinkL}{poppink!20}
\colorlet{popblueL}{popblue!20}
\colorlet{popgoldL}{popgold!20}
\colorlet{popredL}{popred!20}
\colorlet{popgrayL}{popgray!20}
% Teintes claires pour fonds de tableaux / zones
\colorlet{popblueL}{popblue!10!white}
\colorlet{poporangeL}{poporange!14!white}
\colorlet{popgreenL}{popgreen!12!white}
\colorlet{poppurpleL}{poppurple!12!white}
\colorlet{poppinkL}{poppink!12!white}
\colorlet{popgoldL}{popgold!14!white}

\pagecolor{popcream}

% ============================================================
% Typographie
% ============================================================
\defaultfontfeatures{Ligatures=TeX}
\setmainfont{PlusJakartaSans-Medium.ttf}[Path=../../../fonts/,BoldFont=PlusJakartaSans-Bold.ttf]
% Symboles, API et emojis absents de la police principale : repli sur DejaVu Sans (caractères actifs)
\newfontface\symfont{DejaVuSans.ttf}[Path=../../../fonts/]
\makeatletter
\newcommand\symactive[1]{\catcode#1=13 \begingroup\lccode`\~=#1 \lowercase{\endgroup\def~}{{\symfont\char#1}}}
\newcommand\symblank[1]{\catcode#1=13 \begingroup\lccode`\~=#1 \lowercase{\endgroup\def~}{}}
\newcommand\symhyph[1]{\catcode#1=13 \begingroup\lccode`\~=#1 \lowercase{\endgroup\def~}{\mbox{-}}}
\newcount\symc
\newcommand\symrange[2]{\symc=#1 \@whilenum\symc<#2 \do{\expandafter\symactive\expandafter{\the\symc}\advance\symc by 1 }}
\newcommand\symblankrange[2]{\symc=#1 \@whilenum\symc<#2 \do{\expandafter\symblank\expandafter{\the\symc}\advance\symc by 1 }}
\makeatother
\symrange{"0250}{"02FF}\symactive{"2713}\symactive{"2714}\symactive{"2717}\symactive{"2718}\symactive{"2610}\symactive{"2611}\symactive{"2612}
\symhyph{"2011}\symhyph{"2E1A}\symblankrange{"1F300}{"1FAFF}\symblank{"FE0F}
% Maths en sans-serif (Fira Math) avec chiffres et lettres droites de Plus
% Jakarta, pour que les formules se fondent dans le texte.
\usepackage[math-style=ISO,bold-style=ISO]{unicode-math}
\setmathfont{FiraMath-Regular.otf}[Path=../../../fonts/]
\setmathfont{PlusJakartaSans-Medium.ttf}[Path=../../../fonts/,range={up/num,up/{Latin,latin},\mathup}]
\usepackage{icomma} % virgule décimale sans espace en mode math
% Caractères mathématiques Unicode absents des polices de texte (Plus Jakarta,
% Archivo Black) : rendus par leur équivalent mathématique, en texte comme en
% formule — sinon ils s'affichent en carré vide (ex. « Arithmétique dans ℤ »).
\usepackage{newunicodechar}
\newunicodechar{ℝ}{\ensuremath{\mathbb{R}}}
\newunicodechar{ℤ}{\ensuremath{\mathbb{Z}}}
\newunicodechar{ℂ}{\ensuremath{\mathbb{C}}}
\newunicodechar{ℕ}{\ensuremath{\mathbb{N}}}
\newunicodechar{ℚ}{\ensuremath{\mathbb{Q}}}
\newunicodechar{𝔻}{\ensuremath{\mathbb{D}}}
\newunicodechar{α}{\ensuremath{\alpha}}
\newunicodechar{β}{\ensuremath{\beta}}
\newunicodechar{γ}{\ensuremath{\gamma}}
\newunicodechar{δ}{\ensuremath{\delta}}
\newunicodechar{ε}{\ensuremath{\varepsilon}}
\newunicodechar{θ}{\ensuremath{\theta}}
\newunicodechar{λ}{\ensuremath{\lambda}}
\newunicodechar{μ}{\ensuremath{\mu}}
\newunicodechar{σ}{\ensuremath{\sigma}}
\newunicodechar{τ}{\ensuremath{\tau}}
\newunicodechar{φ}{\ensuremath{\varphi}}
\newunicodechar{ψ}{\ensuremath{\psi}}
\newunicodechar{ω}{\ensuremath{\omega}}
\newunicodechar{Σ}{\ensuremath{\Sigma}}
% majuscules produites par \MakeUppercase dans les titres (α-aminés -> Α)
\newunicodechar{Α}{\ensuremath{\alpha}}
\newunicodechar{Β}{\ensuremath{\beta}}
\newunicodechar{Γ}{\ensuremath{\Gamma}}
\newunicodechar{Δ}{\ensuremath{\Delta}}
\newunicodechar{Ε}{\ensuremath{\varepsilon}}
\newunicodechar{Θ}{\ensuremath{\Theta}}
\newunicodechar{Λ}{\ensuremath{\Lambda}}
\newunicodechar{Μ}{\ensuremath{\mu}}
\newunicodechar{Τ}{\ensuremath{\tau}}
\newunicodechar{Φ}{\ensuremath{\Phi}}
\newunicodechar{Ψ}{\ensuremath{\Psi}}
\newunicodechar{Ω}{\ensuremath{\Omega}}
\newunicodechar{↦}{\ensuremath{\mapsto}}
\newunicodechar{∈}{\ensuremath{\in}}
\newunicodechar{∉}{\ensuremath{\notin}}
\newunicodechar{∧}{\ensuremath{\wedge}}
\newunicodechar{⇔}{\ensuremath{\Leftrightarrow}}
\newunicodechar{⟺}{\ensuremath{\Longleftrightarrow}}
\newunicodechar{⇒}{\ensuremath{\Rightarrow}}
\newunicodechar{⟹}{\ensuremath{\Longrightarrow}}
\newunicodechar{∘}{\ensuremath{\circ}}
\newunicodechar{∪}{\ensuremath{\cup}}
\newunicodechar{∩}{\ensuremath{\cap}}
\newunicodechar{≡}{\ensuremath{\equiv}}
\newunicodechar{⊂}{\ensuremath{\subset}}
\newunicodechar{⊥}{\ensuremath{\perp}}
\newunicodechar{∃}{\ensuremath{\exists}}
\newunicodechar{∀}{\ensuremath{\forall}}
\newunicodechar{∛}{\ensuremath{\sqrt[3]{\,}}}
\newunicodechar{∜}{\ensuremath{\sqrt[4]{\,}}}
\newunicodechar{⃗}{\ensuremath{{}^{\to}}}
\newunicodechar{ⁿ}{\ifmmode^{n}\else\textsuperscript{\ensuremath{n}}\fi}
\newunicodechar{ˣ}{\ifmmode^{x}\else\textsuperscript{\ensuremath{x}}\fi}
\newunicodechar{ᵇ}{\ifmmode^{b}\else\textsuperscript{\ensuremath{b}}\fi}
\newunicodechar{ʳ}{\ifmmode^{r}\else\textsuperscript{\ensuremath{r}}\fi}
\newunicodechar{ᵘ}{\ifmmode^{u}\else\textsuperscript{\ensuremath{u}}\fi}
\newunicodechar{ʸ}{\ifmmode^{y}\else\textsuperscript{\ensuremath{y}}\fi}
\newunicodechar{ᵃ}{\ifmmode^{a}\else\textsuperscript{\ensuremath{a}}\fi}
\newunicodechar{ᵉ}{\ifmmode^{e}\else\textsuperscript{\ensuremath{e}}\fi}
\newunicodechar{ᵏ}{\ifmmode^{k}\else\textsuperscript{\ensuremath{k}}\fi}
\newunicodechar{ᵖ}{\ifmmode^{p}\else\textsuperscript{\ensuremath{p}}\fi}
\newunicodechar{ᴺ}{\ifmmode^{N}\else\textsuperscript{\ensuremath{N}}\fi}
\newunicodechar{ᶜ}{\ifmmode^{c}\else\textsuperscript{\ensuremath{c}}\fi}
\newunicodechar{ʰ}{\ifmmode^{h}\else\textsuperscript{\ensuremath{h}}\fi}
\newunicodechar{ᵠ}{\ifmmode^{q}\else\textsuperscript{\ensuremath{q}}\fi}
\newunicodechar{ₙ}{\ifmmode_{n}\else\textsubscript{\ensuremath{n}}\fi}
\newunicodechar{ₐ}{\ifmmode_{a}\else\textsubscript{\ensuremath{a}}\fi}
\newunicodechar{ᵢ}{\ifmmode_{i}\else\textsubscript{\ensuremath{i}}\fi}
\newunicodechar{ᵣ}{\ifmmode_{r}\else\textsubscript{\ensuremath{r}}\fi}
\newunicodechar{ₖ}{\ifmmode_{k}\else\textsubscript{\ensuremath{k}}\fi}
\newunicodechar{ᵦ}{\ifmmode_{\beta}\else\textsubscript{\ensuremath{\beta}}\fi}
\newunicodechar{ₓ}{\ifmmode_{x}\else\textsubscript{\ensuremath{x}}\fi}
\newfontfamily\popdisplay{ArchivoBlack-Regular.ttf}[Path=../../../fonts/]
\newfontfamily\poplogo{SpaceGrotesk-Bold.ttf}[Path=../../../fonts/]
\newfontfamily\popsemi{PlusJakartaSans-SemiBold.ttf}[Path=../../../fonts/]
\newfontfamily\popxbold{PlusJakartaSans-ExtraBold.ttf}[Path=../../../fonts/]
\newfontfamily\popxboldls{PlusJakartaSans-ExtraBold.ttf}[Path=../../../fonts/,LetterSpace=3.0]

\setlist[enumerate]{leftmargin=1.7em,itemsep=5pt,topsep=4pt}
\setlist[itemize]{leftmargin=1.7em,itemsep=4pt,topsep=4pt,label={\color{poporange}\textbullet}}
\setlength{\parindent}{0pt}
\setlength{\parskip}{5pt}
\renewcommand{\arraystretch}{1.25}

% pgfplots : style Pop par défaut
\pgfplotsset{
  every axis/.append style={axis line style={popink,line width=1pt},tick style={popink},
    grid style={popline,line width=0.5pt},label style={font=\small},tick label style={font=\footnotesize}},
  popcurve/.style={poporange,line width=1.8pt,no markers,smooth},
}
% Même style disponible dans un \draw TikZ simple (hors pgfplots)
\tikzset{popcurve/.style={poporange,line width=1.8pt}}
\tikzset{popgrid/.style={popline,very thin}}
\colorlet{popcurve}{poporange} % le nom du style est parfois utilisé comme couleur

% ============================================================
% Pill / squiggle
% ============================================================
\newcommand{\poppill}[3]{%
  \tikz[baseline=-0.55ex]\node[draw=popink,line width=1.2pt,rounded corners=2.6pt,%
    fill=#2,inner xsep=6.5pt,inner ysep=3pt]{%
    \popxboldls\fontsize{7}{7}\selectfont\color{#3}\MakeUppercase{#1}};%
}
\newcommand{\popsquiggle}[1]{%
  \tikz[baseline=-0.4ex]{
    \draw[#1,line width=2.6pt,line cap=round]
      plot [smooth] coordinates {(0,0.09) (0.34,0.01) (0.68,0.21) (1.02,0.01) (1.36,0.10)};
  }%
}
% Mot-clé surligné (marqueur orange sous le texte)
\newcommand{\cle}[1]{{\popxbold\color{popdark}#1}}

% ============================================================
% En-tête / pied
% ============================================================
\pagestyle{fancy}
\fancyhf{}
\renewcommand{\headrulewidth}{0pt}
\renewcommand{\footrulewidth}{0pt}
\lhead{\raisebox{-0.15ex}{\poplogo\fontsize{13}{13}\selectfont\color{popink}OnBuch\color{poporange}+}}
\rhead{\poppill{\DOCMATIERE\ \textbullet\ \DOCNIVEAU\ \textbullet\ Leçon \DOCLECON}{white}{popink}}
\lfoot{\popxbold\fontsize{7.6}{7.6}\selectfont\color{popmuted}\MakeUppercase{Cours signé OnBuch+}\ \textbullet\ {\popsemi\DOCTITRE}}
\rfoot{\tikz[baseline=-0.6ex]\node[rounded corners=8.5pt,fill=popink,minimum height=17pt,minimum width=17pt,inner xsep=5pt,inner sep=0]{\popxbold\fontsize{8}{8}\selectfont\color{popcream}\thepage\,/\,\pageref*{LastPage}};}

% ============================================================
% Titres
% ============================================================
\newcommand{\chaptitre}[2]{%
  \begin{tcolorbox}[enhanced,colback=poporange,colframe=popink,boxrule=2.6pt,arc=11pt,
      drop shadow=popink,left=15pt,right=15pt,top=12pt,bottom=11pt,before skip=3pt,after skip=18pt]
    \poppill{#2}{popink}{popcream}\par\vspace{9pt}
    {\raggedright\hyphenpenalty=10000\exhyphenpenalty=10000\popdisplay\fontsize{22}{24}\selectfont\color{popcream}\MakeUppercase{#1}\par}
  \end{tcolorbox}
}
% Section numérotée : pastille orange avec le numéro + titre Archivo
\newcounter{popsec}
\newcommand{\coursec}[1]{%
  \stepcounter{popsec}\setcounter{popsub}{0}%
  \par\vspace{16pt}\noindent
  \begin{minipage}{\linewidth}
  \tikz[baseline=-0.7ex]\node[draw=popink,line width=1.6pt,fill=poporange,rounded corners=4pt,minimum size=22pt,inner sep=0]{\popdisplay\fontsize{12}{12}\selectfont\color{popcream}\thepopsec};%
  \hspace{8pt}{\popdisplay\fontsize{15}{17}\selectfont\color{popink}\MakeUppercase{#1}}\par
  \vspace{4pt}\noindent{\color{poporange}\rule{36pt}{3.6pt}}
  \end{minipage}\par\nopagebreak\vspace{8pt}
}
\newcounter{popsub}
\newcommand{\courssub}[1]{%
  \stepcounter{popsub}%
  \par\vspace{9pt}\noindent{\popxbold\fontsize{11.5}{13}\selectfont\color{popink}{\color{poporange}\thepopsec.\thepopsub}\hspace{6pt}#1}\par\nopagebreak\vspace{4pt}
}

% ============================================================
% Boîtes pédagogiques — même recette : fond blanc, bordure épaisse colorée,
% ombre dure encre, badge plein. Titre optionnel [#1].
% ============================================================
\tcbset{popbase/.style={breakable,enhanced,colback=white,boxrule=2.3pt,arc=9pt,
  shadow={3pt}{-3pt}{0pt}{popink},left=14pt,right=14pt,top=11pt,bottom=11pt,before skip=11pt,after skip=15pt,
  fontupper=\popsemi\fontsize{10.6}{14.5}\selectfont\color{popink},
  fontlower=\popsemi\fontsize{10.6}{14.5}\selectfont\color{popink}}}
% Coche dessinée (le glyphe ✓ n'existe pas dans les polices du document)
\newcommand{\popcheck}[1]{\tikz[baseline=-0.5ex]\draw[#1,line width=1.6pt,line cap=round,line join=round](-0.13,0.0)--(-0.04,-0.09)--(0.13,0.1);}
\providecommand{\coche}{\popcheck{popgreen}}
\providecommand{\resultat}[1]{\par\smallskip\noindent\fcolorbox{popgreen}{popgreenL}{\parbox{\dimexpr\linewidth-2\fboxsep-2\fboxrule\relax}{\textbf{Résultat.} #1}}\par\smallskip}
\newcommand{\popboxhead}[3]{\poppill{#1}{#2}{white}\ifx\relax#3\relax\else\hspace{6pt}{\popxbold\fontsize{10.6}{12}\selectfont\color{popink}#3}\fi\par\vspace{7pt}}

\newtcolorbox{aretenir}[1][]{popbase,colframe=popgreen,before upper={\popboxhead{À retenir}{popgreen}{#1}}}
% Texte en anglais (document de lecture, dialogue, extrait) : cadre
% d'étude. Un titre optionnel (source, auteur) s'affiche dans l'en-tête.
\newtcolorbox{textebox}[1][]{popbase,colframe=popblue,before upper={\popboxhead{Text}{popblue}{#1}\begin{otherlanguage}{english}},after upper={\end{otherlanguage}}}
% Marques ✓ / ✗ (forme correcte / fautive) souvent utilisées par les modèles
\providecommand{\cross}{\textcolor{poppink}{\ensuremath{\times}}}
\providecommand{\xmark}{\cross}\providecommand{\crossmark}{\cross}\providecommand{\cmark}{\ensuremath{\checkmark}}
% Ligne de réponse / justification à compléter (inventée par les modèles)
\providecommand{\justification}[1]{\par\noindent\textit{Justification~:} \dotfill\par}
% \textipa (package tipa non chargé) : la police ne couvre pas l'API -> simple passage
\providecommand{\textipa}[1]{#1}
% Soulignement (ulem non chargé) : \uline, \ul
\providecommand{\uline}[1]{\underline{#1}}\providecommand{\ul}[1]{\underline{#1}}
% \glossaire{\item ...} inventée par les modèles : liste de glossaire
\providecommand{\glossaire}[1]{\begin{itemize}#1\end{itemize}}
% \alert (beamer) inventée par les modèles : texte mis en évidence
\providecommand{\alert}[1]{\textbf{\textcolor{poppink}{#1}}}
% variantes ASCII d'ordinaux écrites en macro
\providecommand{\iere}{\textsuperscript{re}}\providecommand{\ieme}{\textsuperscript{e}}\providecommand{\ere}{\textsuperscript{re}}\providecommand{\eme}{\textsuperscript{e}}\providecommand{\er}{\textsuperscript{er}}\providecommand{\ie}{\textsuperscript{e}}
% Ordinaux français écrits en macro par les modèles : 1\ière, 3\ième
\newcommand{\ière}{\textsuperscript{re}}\newcommand{\ième}{\textsuperscript{e}}
% \exemple (inventée par les modèles) : étiquette « Exemple : »
\providecommand{\exemple}{\textbf{Exemple~:}\ }
% \square utilisable aussi hors mode math (cases à cocher)
\AtBeginDocument{\let\squareMath\square\renewcommand{\square}{\ensuremath{\squareMath}}}
% Mot ou phrase en anglais à l'intérieur d'un paragraphe français
\newcommand{\eng}[1]{\ifmmode\text{\textit{#1}}\else\begingroup\selectlanguage{english}\itshape #1\endgroup\fi}
\let\esp\eng % alias toléré
% flèches d'intonation inventées par les modèles
\providecommand{\textsearrow}{}\renewcommand{\textsearrow}{\ensuremath{\searrow}}\providecommand{\textnearrow}{}\renewcommand{\textnearrow}{\ensuremath{\nearrow}}
% \fr{..} : mot français (inventé par les modèles)
\providecommand{\fr}[1]{\textit{#1}}
\newtcolorbox{exemplebox}[1][]{popbase,colframe=poporange,before upper={\popboxhead{Exemple}{poporange}{#1}}}
\newtcolorbox{definition}[1][]{popbase,colframe=popblue,before upper={\popboxhead{Définition}{popblue}{#1}}}
\newtcolorbox{propriete}[1][]{popbase,colframe=poppurple,before upper={\popboxhead{Propriété}{poppurple}{#1}}}
\newtcolorbox{methode}[1][]{popbase,colframe=popgold,before upper={\popboxhead{Méthode}{popgold}{#1}}}
\newtcolorbox{attention}[1][]{popbase,colframe=poppink,before upper={\popboxhead{Attention}{poppink}{#1}}}
\newtcolorbox{experience}[1][]{popbase,colframe=popblue,colback=popblueL,before upper={\popboxhead{Expérience}{popblue}{#1}}}
% « Le savais-tu ? » : carte violette pleine (contexte, culture, Cameroun)
\newtcolorbox{savaistu}[1][]{popbase,colback=poppurple,colframe=popink,
  fontupper=\popsemi\fontsize{10.4}{14.2}\selectfont\color{white},
  before upper={\poppill{Le savais-tu\,?}{white}{poppurple}\ifx\relax#1\relax\else\hspace{6pt}{\popxbold\color{white}#1}\fi\par\vspace{7pt}}}
% Exercice résolu : énoncé en haut, \tcblower puis la solution sur fond crème
\newcounter{popexo}\newcounter{popexr}
\newtcolorbox{exoresolu}[1][]{popbase,colframe=poporange,colbacklower=popcream,
  segmentation style={popink,line width=1.2pt,dashed},
  before upper={\stepcounter{popexr}\popboxhead{Exercice résolu \thepopexr}{poporange}{#1}},
  before lower={\poppill{Solution}{popink}{popcream}\par\vspace{6pt}}}
% Exercice (non résolu en place) — la correction est dans la partie Corrigés
\newtcolorbox{exercice}[1][]{popbase,colframe=popink,
  before upper={\stepcounter{popexo}\popboxhead{Exercice \thepopexo}{popink}{#1}}}
% Corrigé
\newtcolorbox{corrige}[1][]{popbase,colframe=popgreen,colback=popgreenL,
  before upper={\popboxhead{Corrigé}{popgreen}{#1}}}
% Objectifs / prérequis / bilan
\newtcolorbox{objectifs}{popbase,colframe=popink,colback=poporangeL,
  before upper={\popboxhead{Objectifs de la leçon}{popink}{}}}
\newtcolorbox{prerequis}{popbase,colframe=popink,colback=popblueL,
  before upper={\popboxhead{Prérequis}{popblue}{}}}
\newtcolorbox{bilan}{popbase,colframe=popink,colback=white,boxrule=2.8pt,
  before upper={\popboxhead{Fiche bilan}{poporange}{L'essentiel à connaître pour l'examen}}}
% Figure encadrée avec légende
\newtcolorbox{popfigure}[1][]{enhanced,breakable=false,colback=white,colframe=popline,boxrule=1.4pt,arc=8pt,
  left=10pt,right=10pt,top=10pt,bottom=8pt,before skip=10pt,after skip=12pt,halign=center,
  lower separated=false,halign lower=center,
  fontlower=\popsemi\fontsize{9}{11.5}\selectfont\color{popmuted},
  #1}
\newcommand{\legende}[1]{\par\vspace{6pt}{\popsemi\fontsize{9}{11.5}\selectfont\color{popmuted}#1}}

% ============================================================
% Signature OnBuch+
% ============================================================
\newcommand{\poplogoinline}[1]{{\poplogo\fontsize{#1}{#1}\selectfont\color{popink}OnBuch\color{poporange}+}}
\newcommand{\signatureonbuch}{%
  \begin{tcolorbox}[enhanced,colback=popink,colframe=popink,boxrule=2.2pt,arc=10pt,
      drop shadow=poporange,left=14pt,right=14pt,top=10pt,bottom=10pt,before skip=0pt,after skip=0pt]
    \tikz[baseline=-0.7ex]\node[circle,fill=popgreen,draw=popcream,line width=1.3pt,minimum size=22pt,inner sep=0]{\popcheck{white}};%
    \hspace{9pt}\begin{minipage}[c]{0.62\linewidth}
      {\popxbold\fontsize{12}{13}\selectfont\color{popcream}Signé\ \textbullet\ {\poplogo OnBuch\color{poporange}+}}\par\vspace{2pt}
      {\raggedright\popsemi\fontsize{8}{10.5}\selectfont\color{popline}\MakeUppercase{Équipe pédagogique OnBuch+}\\\MakeUppercase{Conforme au programme officiel MINESEC}\par}
    \end{minipage}\hfill
    {\poplogo\fontsize{22}{22}\selectfont\color{popcream}OnBuch\color{poporange}+}
  \end{tcolorbox}%
}

% ============================================================
% Page de garde
% #1 titre · #2 matière · #3 niveau · #4 module · #5 n° leçon · #6 durée ·
% #7 description / objectifs (texte ou itemize)
% ============================================================
\newcommand{\pagegarde}[7]{%
  \thispagestyle{empty}
  \newgeometry{a4paper,margin=1.7cm,top=1.6cm,bottom=1.4cm}%
  \begin{tikzpicture}[remember picture,overlay]
    \fill[poporange!18] ([xshift=-3.2cm,yshift=-3.6cm]current page.north east) circle (4.6cm);
    \fill[poppurple!14] ([xshift=2.4cm,yshift=7.6cm]current page.south west) circle (3.8cm);
  \end{tikzpicture}%
  {\poplogo\fontsize{30}{30}\selectfont\color{popink}OnBuch\color{poporange}+}\hfill
  \raisebox{0.6ex}{\poppill{Cours signé \textbullet\ Premium}{popink}{popcream}}\par
  \vspace{1pt}\popsquiggle{poppurple}\par
  \vspace{8pt}
  \begin{tcolorbox}[enhanced,colback=poporange,colframe=popink,boxrule=2.8pt,arc=13pt,
      drop shadow=popink,left=18pt,right=18pt,top=12pt,bottom=13pt,after skip=10pt]
    \poppill{#2 \textbullet\ #3}{popink}{popcream}\hspace{5pt}\poppill{#4}{white}{popink}\par\vspace{8pt}
    {\popdisplay\fontsize{14}{14}\selectfont\color{popink}LEÇON #5}\par\vspace{4pt}
    {\raggedright\hyphenpenalty=10000\exhyphenpenalty=10000\popdisplay\fontsize{25}{27}\selectfont\color{popcream}\MakeUppercase{#1}\par}
    \vspace{8pt}
    {\popxbold\fontsize{9.5}{9.5}\selectfont\color{popink}\MakeUppercase{Durée conseillée : #6}}
  \end{tcolorbox}
  \begin{tcolorbox}[enhanced,colback=white,colframe=popink,boxrule=2.2pt,arc=10pt,
      drop shadow=popink,left=16pt,right=16pt,top=10pt,bottom=10pt,after skip=8pt]
    \poppill{Dans cette leçon}{poppurple}{white}\par\vspace{5pt}
    \popsemi\fontsize{9.3}{12.6}\selectfont\color{popink}#7
  \end{tcolorbox}
  \noindent\begin{minipage}[t]{0.487\linewidth}
    \begin{tcolorbox}[enhanced,colback=popgreen,colframe=popink,boxrule=2.2pt,arc=10pt,
        drop shadow=popink,left=11pt,right=11pt,top=10pt,bottom=10pt,height=3.1cm,valign=center]
      \tcbox[colback=white,colframe=popink,boxrule=1.3pt,arc=5pt,boxsep=0pt,nobeforeafter,
        left=3pt,right=3pt,top=3pt,bottom=3pt,valign=center]{\qrcode[height=1.9cm]{https://play.google.com/store/apps/details?id=com.luvvix.onbuch&hl=fr}}%
      \hspace{8pt}\begin{minipage}[c]{0.5\linewidth}
        {\popdisplay\fontsize{11}{12.5}\selectfont\color{white}\MakeUppercase{Télécharge OnBuch}}\par\vspace{4pt}
        {\raggedright\popsemi\fontsize{8.4}{11}\selectfont\color{white}Gratuit \textbullet\ Google Play\\Tuteur IA, annales, concours, résultats\par}
      \end{minipage}
    \end{tcolorbox}
  \end{minipage}\hfill
  \begin{minipage}[t]{0.487\linewidth}
    \begin{tcolorbox}[enhanced,colback=poppurple,colframe=popink,boxrule=2.2pt,arc=10pt,
        drop shadow=popink,left=11pt,right=11pt,top=10pt,bottom=10pt,height=3.1cm,valign=center]
      \tikz[baseline=-0.7ex]\node[circle,fill=white,draw=popink,line width=1.3pt,minimum size=1.6cm,inner sep=0]{\popdisplay\fontsize{24}{24}\selectfont\color{poppurple}+};%
      \hspace{8pt}\begin{minipage}[c]{0.6\linewidth}
        {\popdisplay\fontsize{11}{12.5}\selectfont\color{white}\MakeUppercase{Passe à OnBuch+}}\par\vspace{4pt}
        {\raggedright\popsemi\fontsize{8.4}{11}\selectfont\color{white}Tous les cours signés, Tuteur IA illimité, fiches PDF — onglet OnBuch+ de l'app\par}
      \end{minipage}
    \end{tcolorbox}
  \end{minipage}
  \vspace{8pt}
  \popcontenu
  \vfill
  \signatureonbuch
  \restoregeometry
  \newpage
}

% Rangée « Ce cours contient » (page de garde)
\newcommand{\poptile}[3]{% #1 couleur #2 gros texte #3 légende
  \begin{tcolorbox}[enhanced,colback=#1,colframe=popink,boxrule=2pt,arc=9pt,shadow={2.5pt}{-2.5pt}{0pt}{popink},
      left=6pt,right=6pt,top=9pt,bottom=9pt,halign=center,valign=center,height=1.75cm,width=\linewidth,nobeforeafter]
    {\popdisplay\fontsize{12.5}{13}\selectfont\color{white}\MakeUppercase{#2}}\par\vspace{3pt}
    {\centering\popsemi\fontsize{7.8}{9.5}\selectfont\color{white}#3\par}
  \end{tcolorbox}}
\newcommand{\popcontenu}{%
  \par\noindent{\popxboldls\fontsize{7.5}{7.5}\selectfont\color{popmuted}CE COURS SIGNÉ CONTIENT}\par\vspace{7pt}
  \noindent\begin{minipage}[t]{0.235\linewidth}\poptile{popblue}{Cours}{complet, illustré, pas à pas}\end{minipage}\hfill
  \begin{minipage}[t]{0.235\linewidth}\poptile{poporange}{Méthodes}{et exercices résolus}\end{minipage}\hfill
  \begin{minipage}[t]{0.235\linewidth}\poptile{poppink}{Exercices}{type examen + corrigés}\end{minipage}\hfill
  \begin{minipage}[t]{0.235\linewidth}\poptile{popgreen}{Fiche bilan}{l'essentiel à retenir}\end{minipage}\par}

% Sommaire
\newtcolorbox{sommaire}{enhanced,colback=white,colframe=popink,boxrule=2.2pt,arc=10pt,
  drop shadow=popink,left=18pt,right=18pt,top=13pt,bottom=13pt,before skip=6pt,after skip=20pt,
  before upper={\poppill{Sommaire}{poppurple}{white}\par\vspace{9pt}\popsemi\fontsize{10.6}{14.5}\selectfont\color{popink}}}

% Signature de fin de cours — poussée tout en bas de la dernière page
\newcommand{\findecours}{%
  \par\vfill\nopagebreak
  \begin{center}\popsquiggle{poporange}\end{center}\vspace{4pt}
  \signatureonbuch
}
\AtBeginDocument{\def\llbracket{\mathopen{[\mkern-3.5mu[}}\def\rrbracket{\mathclose{]\mkern-3.5mu]}}} % intervalles d'entiers (glyphe ⟦ absent de la police)
% Glyphes absents de Fira Math : redéfinis à partir de glyphes disponibles
\AtBeginDocument{%
  \def\setminus{\mathbin{\backslash}}\let\smallsetminus\setminus
  \def\vdots{\mathord{\vbox{\baselineskip3.2pt\lineskiplimit0pt\kern4pt\hbox{.}\hbox{.}\hbox{.}}}}%
  \def\ddots{\mathinner{\mkern1mu\raise6.5pt\hbox{.}\mkern2mu\raise3.6pt\hbox{.}\mkern2mu\raise0.7pt\hbox{.}\mkern1mu}}%
  \def\triangle{\mathord{\scalebox{0.9}{$\symup{\Delta}$}}}%
  \def\nabla{\mathord{\raisebox{\depth}{\scalebox{1}[-1]{$\symup{\Delta}$}}}}%
  \def\checkmark{\popcheck{popdark}}%
  \def\star{\mathbin{\ast}}\def\bigstar{\mathord{\ast}}%
  \def\diamond{\mathbin{\scalebox{0.6}{$\Diamond$}}}%
  \def\bigcup{\mathop{\mathchoice{\raisebox{-0.3em}{\scalebox{1.7}{$\cup$}}}{\scalebox{1.25}{$\cup$}}{\cup}{\cup}}\displaylimits}%
  \def\bigcap{\mathop{\mathchoice{\raisebox{-0.3em}{\scalebox{1.7}{$\cap$}}}{\scalebox{1.25}{$\cap$}}{\cap}{\cap}}\displaylimits}%
  \def\lesseqqgtr{\mathrel{\lessgtr}}\def\bigoplus{\mathop{\oplus}\displaylimits}%
}
\providecommand{\ssi}{si et seulement si} % abréviation fréquente
\AtBeginDocument{\providecommand{\wideparen}{\overparen}} % arc de cercle

%% FIGFIX-BEGIN v3 — correctifs globaux des figures (docs/onbuch-generation/tools/figfix)
\usepackage{adjustbox}\usepackage{etoolbox}
% toute figure TikZ/pgfplots plus large que la ligne est réduite à la largeur disponible
\BeforeBeginEnvironment{tikzpicture}{\begin{adjustbox}{max width=\linewidth}}
\AfterEndEnvironment{tikzpicture}{\end{adjustbox}}
% légendes pgfplots lisibles par-dessus les courbes
\pgfplotsset{every axis/.append style={legend style={fill=white,fill opacity=0.92,text opacity=1,draw=black!15,inner sep=3pt}}}
% étiquettes d'arêtes (syntaxe edge["..."]) avec fond blanc
\tikzset{every edge quotes/.append style={fill=white,inner sep=1.2pt}}
% bulles de cartes mentales : texte à la ligne dans la bulle, bulle un peu plus grande
\tikzset{every concept/.append style={text width=2.0cm,align=center,minimum size=2.3cm,inner sep=2pt}}
%% FIGFIX-END
\begin{document}

\pagegarde{Lesson 34 — Grammar: More difficult structures with adjectives: How + adj, the superlative, too, to, etc}{Anglais}{Tle A}{Chapter 3 : Environment, Well-being and Health}{EN34}{2 h}{Cette leçon de grammaire étudie les structures complexes construites autour des adjectifs : les questions en « How + adjectif », la formation et l'emploi du superlatif, ainsi que les constructions « too + adjectif + to-infinitive » et « adjectif + enough + to-infinitive ». Elle vise à corriger les erreurs fréquentes des francophones et à utiliser ces structures dans des phrases et de courts textes liés à l'environnement, au bien-être et à la santé.
\vspace{4pt}
\begin{multicols}{2}\small
\begin{itemize}
\item À la fin de cette leçon, je sais poser des questions avec « How + adjectif » (How old, How tall, How far, How long...) et y répondre correctement.
\item À la fin de cette leçon, je sais former le superlatif des adjectifs courts (-est) et longs (the most) sans erreur.
\item À la fin de cette leçon, je sais utiliser les superlatifs irréguliers (the best, the worst, the farthest/furthest).
\item À la fin de cette leçon, je sais employer la structure « too + adjectif (+ for + nom) + to-infinitive » pour exprimer un excès empêchant une action.
\item À la fin de cette leçon, je sais employer la structure « adjectif + enough (+ for + nom) + to-infinitive » pour exprimer une suffisance.
\item À la fin de cette leçon, je sais transformer une phrase avec « too » en phrase avec « enough » (et inversement) en conservant le sens.
\end{itemize}\end{multicols}}

\begin{sommaire}
\begin{multicols}{2}
\begin{enumerate}
\item Observation : découvrir les structures en contexte
\item How + adjectif : les questions de mesure
\item Le superlatif : formation et emplois
\item Too...to et enough...to : exprimer la limite
\item Pièges des francophones et entraînement ciblé
\item Activité d'intégration
\item Méthodes et exercices résolus
\item Exercices
\item Corrigés des exercices
\item Fiche bilan
\end{enumerate}
\end{multicols}
\end{sommaire}

\begin{objectifs}
\begin{itemize}
\item À la fin de cette leçon, je sais poser des questions avec « How + adjectif » (How old, How tall, How far, How long...) et y répondre correctement.
\item À la fin de cette leçon, je sais former le superlatif des adjectifs courts (-est) et longs (the most) sans erreur.
\item À la fin de cette leçon, je sais utiliser les superlatifs irréguliers (the best, the worst, the farthest/furthest).
\item À la fin de cette leçon, je sais employer la structure « too + adjectif (+ for + nom) + to-infinitive » pour exprimer un excès empêchant une action.
\item À la fin de cette leçon, je sais employer la structure « adjectif + enough (+ for + nom) + to-infinitive » pour exprimer une suffisance.
\item À la fin de cette leçon, je sais transformer une phrase avec « too » en phrase avec « enough » (et inversement) en conservant le sens.
\item À la fin de cette leçon, je sais éviter les erreurs typiques des francophones (She has 20 years, the more big, too much small, enough old...).
\item À la fin de cette leçon, je sais utiliser ces structures dans un court paragraphe descriptif sur la santé ou l'environnement.
\end{itemize}
\end{objectifs}

\begin{prerequis}
\begin{itemize}
\item Le comparatif de supériorité (taller than, more beautiful than) vu en Première
\item La formation des adjectifs et leur place devant le nom (a clean hospital)
\item Le présent simple et la question avec les auxiliaires do/does
\item Les adverbes de quantité much, many, a lot of
\item L'infinitif avec to après un verbe (I want to go)
\end{itemize}
\end{prerequis}

\newpage

\coursec{Observation : découvrir les structures en contexte}

\courssub{Situation d'accroche : trois petites scènes de la vie quotidienne}

Pendant la campagne de vaccination à l'hôpital de district de Bamenda, une infirmière anglophone demande à un élève de Terminale : \eng{“How old is your grandmother?”} Il répond fièrement : \eng{“She has 70 years.”} L'infirmière sourit poliment, mais la phrase est incorrecte : c'est l'erreur typique du francophone, qui traduit mot à mot « elle a 70 ans ». Plus tard, le même élève veut expliquer que la salle d'attente est « trop petite pour accueillir tout le monde » et hésite entre \eng{“too small to welcome”} et \eng{“enough small”} — la première est bonne, la seconde n'existe pas. Enfin, le soir, au journal télévisé, il entend : \eng{“Cameroon is one of the most beautiful countries in Africa.”} Le Cameroun est l'un des plus beaux pays d'Afrique.

Ces trois situations quotidiennes — \cle{décrire} une réalité mesurable, \cle{comparer} en situant un élément au sommet d'un groupe, \cle{exprimer une limite} — exigent de maîtriser les structures difficiles construites avec les adjectifs : \cle{How + adjectif}, \cle{le superlatif}, \cle{too...to} et \cle{enough...to}.

\textbf{Problématique de la leçon.} Comment l'anglais exprime-t-il la mesure (âge, distance, taille, prix), la supériorité absolue et la limite, et pourquoi la traduction mot à mot du français conduit-elle si souvent à l'erreur ?

\begin{methode}[Observer une structure grammaticale : la démarche en cinq étapes]
Avant d'apprendre une règle par cœur, le bon élève l'observe d'abord en contexte. Voici la démarche que nous allons suivre ensemble :
\begin{enumerate}
\item \textbf{Lire} l'exemple dans son contexte, sans dictionnaire.
\item \textbf{Deviner le sens} de la structure grâce aux indices du texte (situation, mots voisins, logique).
\item \textbf{Analyser la forme} : ordre des mots, place de l'adjectif, mot qui accompagne l'adjectif.
\item \textbf{Formuler une règle provisoire} avec ses propres mots (une hypothèse).
\item \textbf{Vérifier} l'hypothèse avec la règle officielle du cours.
\end{enumerate}
Cette démarche, dite \emph{inductive}, est celle de l'approche par les compétences : vous découvrez, vous construisez, puis vous fixez le savoir.
\end{methode}

\courssub{Lecture du texte support : The New Health Centre}

Lisons maintenant un court article de presse locale. Pendant la lecture, repérez mentalement tout ce qui tourne autour des adjectifs.

\begin{textebox}[The New Health Centre in Buea — The Mountain Voice, local newspaper]
The new health centre in Buea is the cleanest in the region. How far is it from your village? It is only two kilometres away, so it is near enough for patients to walk there. But the waiting room is too small to seat everybody, and some nurses are too tired to work night shifts. Doctor Ndi is the most dedicated doctor I know. He is also the tallest man in the hospital, and his consultations are the cheapest in town. How old is the building? Only six months! Yet it is already the busiest health centre on the mountain. The director says the centre is not big enough to serve the whole town, so a second building will open next year. “Our patients are too important to wait in the rain,” he adds with a smile.
\end{textebox}

\textbf{Glossaire.}

\begin{center}
\begin{tabularx}{\linewidth}{|l|l|X|}
\hline
\rowcolor{popblueL} \textbf{Mot anglais} & \textbf{Nature} & \textbf{Traduction} \\
\hline
health centre & nom & centre de santé \\
\hline
waiting room & nom & salle d'attente \\
\hline
to seat & verbe & accueillir (asseoir) \\
\hline
night shift & nom & service de nuit \\
\hline
dedicated & adjectif & dévoué \\
\hline
whole & adjectif & entier, tout entier \\
\hline
\end{tabularx}
\end{center}

\textbf{Questions de compréhension.}
\begin{enumerate}
\item \eng{Where is the health centre?} (Où se trouve le centre de santé ?)
\item \eng{Why can patients walk there?} (Pourquoi les patients peuvent-ils s'y rendre à pied ?)
\item \eng{What is the problem with the waiting room?} (Quel est le problème de la salle d'attente ?)
\item \eng{What will happen next year?} (Que se passera-t-il l'année prochaine ?)
\end{enumerate}

\courssub{Repérage guidé : trois gestes, trois couleurs}

Reprenons le texte phrase par phrase et appliquons la consigne de repérage : \textbf{soulignez} les questions en \eng{How + adjectif}, \textbf{encadrez} les superlatifs, \textbf{entourez} les structures avec \eng{too} et \eng{enough}.

\begin{exemplebox}[Le texte annoté]
\begin{itemize}
\item Questions en \eng{How + adjectif} (à souligner) : \eng{How far is it from your village?} — \eng{How old is the building?}
\item Superlatifs (à encadrer) : \eng{the cleanest in the region} — \eng{the most dedicated doctor} — \eng{the tallest man} — \eng{the cheapest in town} — \eng{the busiest health centre}.
\item Limites (à entourer) : \eng{near enough for patients to walk there} — \eng{too small to seat everybody} — \eng{too tired to work night shifts} — \eng{not big enough to serve the whole town} — \eng{too important to wait in the rain}.
\end{itemize}
\end{exemplebox}

\courssub{Hypothèses des élèves : deviner avant d'apprendre}

Travail en binômes. Pour chaque catégorie repérée, formulez une hypothèse de sens et une hypothèse de forme. Voici ce que des élèves de Terminale A proposent généralement :

\begin{itemize}
\item \textbf{Catégorie 1 — How + adjectif.} Hypothèse de sens : « on pose une question sur une mesure ou une quantité ». Hypothèse de forme : « \eng{How} est suivi d'un adjectif, puis vient le verbe, puis le sujet ». Exemple analysé : \eng{How far / is / it?}
\item \textbf{Catégorie 2 — Le superlatif.} Hypothèse de sens : « on dit qu'un élément est le premier, le meilleur, le plus extrême de tout un groupe ». Hypothèse de forme : « il y a toujours \eng{the} devant, et l'adjectif prend une terminaison en \eng{-est} ou est précédé de \eng{most} ».
\item \textbf{Catégorie 3 — Too et enough.} Hypothèse de sens : « on exprime une limite : trop ou assez pour faire quelque chose ». Hypothèse de forme : « \eng{too} se place avant l'adjectif, \eng{enough} se place après l'adjectif, et le verbe qui suit est à l'infinitif avec \eng{to} ».
\end{itemize}

Gardez ces hypothèses en tête : nous les vérifierons dans les sections suivantes de la leçon.

\begin{attention}[Le piège de la traduction mot à mot]
Ne traduisez jamais le français littéralement ! \eng{“How old is she?”} ne se dit \textbf{jamais} \eng{“How many years does she have?”} en anglais correct : l'anglais \emph{est} un âge (\eng{She is 70.}), il ne le « possède » pas. De même, « trop petit pour accueillir » se dit \eng{too small to seat}, jamais \eng{enough small} : \eng{enough} se place \emph{après} l'adjectif (\eng{big enough}), et \eng{too} \emph{avant} (\eng{too small}). Retenez : le français dit « avoir X ans », l'anglais dit \eng{to be X (years old)}.
\end{attention}

\courssub{Mise en commun : classer les exemples en trois catégories}

Après le travail en binômes, la classe met en commun au tableau. Chaque exemple du texte est rangé dans l'une des trois colonnes suivantes :

\begin{center}
\begin{tabularx}{\linewidth}{|X|X|X|}
\hline
\rowcolor{popblueL} \textbf{Questions (How + adj.)} & \textbf{Superlatifs (the ...-est / the most)} & \textbf{Limites (too / enough + to)} \\
\hline
\eng{How far is it?} « À quelle distance est-ce ? » & \eng{the cleanest} « le plus propre » & \eng{too small to seat everybody} « trop petite pour accueillir tout le monde » \\
\hline
\eng{How old is the building?} « Quel âge a le bâtiment ? » & \eng{the most dedicated doctor} « le médecin le plus dévoué » & \eng{too tired to work} « trop fatiguées pour travailler » \\
\hline
 & \eng{the tallest man} « l'homme le plus grand » & \eng{near enough to walk there} « assez proche pour y aller à pied » \\
\hline
 & \eng{the busiest centre} « le centre le plus fréquenté » & \eng{not big enough to serve the town} « pas assez grand pour desservir la ville » \\
\hline
\end{tabularx}
\end{center}

\begin{popfigure}
\begin{tikzpicture}[mindmap, grow cyclic, every node/.style={concept, text=white, align=center},
root concept/.append style={concept color=popblue, font=\bfseries},
level 1/.append style={level distance=3.6cm, sibling angle=120},
level 2/.append style={level distance=2.4cm, sibling angle=45}]
\node[root concept, align=center]{Adjectives:\\difficult structures}
  child[concept color=poporange]{ node[align=center]{How + adjective\\ \eng{How far is it?}}
    child[concept color=poporange!70]{ node {age: \eng{How old?}} }
    child[concept color=poporange!70]{ node {distance: \eng{How far?}} }
  }
  child[concept color=popgreen]{ node[align=center]{Superlative\\ \eng{the cleanest}}
    child[concept color=popgreen!70]{ node {short adj.: \eng{-est}} }
    child[concept color=popgreen!70]{ node {long adj.: \eng{the most}} }
  }
  child[concept color=poppurple]{ node[align=center]{too / enough + to\\ \eng{too small to seat}}
    child[concept color=poppurple!70]{ node {\eng{too} + adj. + \eng{to}} }
    child[concept color=poppurple!70]{ node {adj. + \eng{enough} + \eng{to}} }
  };
\end{tikzpicture}
\legende{Carte mentale des trois structures difficiles construites avec les adjectifs.}
\end{popfigure}

\courssub{Vérification des hypothèses : premières règles provisoires}

Confrontons maintenant les hypothèses des élèves aux faits observés dans le texte.

\begin{aretenir}[Ce que l'observation nous a déjà appris]
\begin{itemize}
\item \eng{How + adjectif} sert à \textbf{questionner sur une mesure} : âge (\eng{How old?}), distance (\eng{How far?}), taille (\eng{How tall?}), longueur (\eng{How long?}), prix (\eng{How much?}). L'ordre est : \eng{How} + adjectif + verbe + sujet.
\item Le \textbf{superlatif} place un élément au sommet d'un groupe : \eng{the} + adjectif court + \eng{-est} (\eng{the cleanest}) ou \eng{the most} + adjectif long (\eng{the most dedicated}). Il est souvent suivi de \eng{in} + groupe (\eng{in the region}).
\item \eng{Too + adjectif + to} + infinitif exprime un \textbf{excès qui empêche} : \eng{too small to seat everybody}. \eng{Adjectif + enough + to} + infinitif exprime une \textbf{suffisance qui permet} : \eng{near enough to walk there}. Attention à la place : \eng{too} avant l'adjectif, \eng{enough} après.
\end{itemize}
Ces règles seront précisées, systématisées et complétées (formes, exceptions) dans les sections suivantes.
\end{aretenir}

\begin{exoresolu}[Classer et expliquer]
Énoncé. Classez les phrases suivantes dans les trois catégories (question de mesure, superlatif, limite) et traduisez-les : 1) \eng{How long is the Wouri bridge?} 2) \eng{Mount Cameroon is the highest peak in West Africa.} 3) \eng{The boy is too young to travel alone.} 4) \eng{This bag is big enough to carry my books.}
\tcblower
Solution détaillée.
\begin{enumerate}
\item \eng{How long is the Wouri bridge?} — \textbf{Question de mesure} (\eng{How} + adjectif \eng{long} + verbe + sujet). Traduction : « Quelle est la longueur du pont de la Wouri ? »
\item \eng{Mount Cameroon is the highest peak in West Africa.} — \textbf{Superlatif} (\eng{the} + adjectif court \eng{high} + \eng{-est}, suivi de \eng{in} + groupe). Traduction : « Le mont Cameroun est le plus haut sommet d'Afrique de l'Ouest. »
\item \eng{The boy is too young to travel alone.} — \textbf{Limite (excès)} : \eng{too} + adjectif \eng{young} + \eng{to} + infinitif. Traduction : « Le garçon est trop jeune pour voyager seul. »
\item \eng{This bag is big enough to carry my books.} — \textbf{Limite (suffisance)} : adjectif \eng{big} + \eng{enough} + \eng{to} + infinitif. Traduction : « Ce sac est assez grand pour porter mes livres. »
\end{enumerate}
\end{exoresolu}

\begin{experience}[À vous de jouer : observation active en binômes]
\begin{enumerate}
\item Relisez le texte à voix haute à deux : un élève lit, l'autre lève la main à chaque structure cible entendue.
\item Trouvez dans le texte un exemple de chaque catégorie que nous n'avons pas encore cité en classe.
\item Posez-vous mutuellement deux questions en \eng{How + adjectif} sur votre école ou votre quartier : \eng{How far is your house from school? How old is your school building?}
\item Notez une hypothèse personnelle sur la différence de sens entre \eng{too small} et \eng{not big enough} : nous la vérifierons dans la section 4.
\end{enumerate}
\end{experience}

\begin{savaistu}[Pourquoi l'anglais « est » son âge ?]
En anglais, l'âge, la faim et la soif se construisent avec le verbe \eng{to be} : \eng{I am twenty}, \eng{I am hungry}, \eng{I am thirsty}. Le français fait l'inverse avec \emph{avoir} : « j'ai vingt ans », « j'ai faim », « j'ai soif ». C'est l'une des différences les plus profondes entre les deux langues, et elle explique à elle seule la moitié des erreurs des francophones dans les questions en \eng{How + adjectif}. Au Cameroun anglophone, à Bamenda ou à Buea, personne ne dira jamais \eng{“I have twenty years”} : on dira \eng{“I am twenty (years old)”}.
\end{savaistu}

\coursec{How + adjectif : les questions de mesure}

Dans la section précédente, nous avons observé un texte sur l'environnement et la santé dans lequel apparaissaient des questions comme \eng{How old is the baobab?} ou \eng{How far is the nearest health centre?}. Nous allons maintenant étudier de manière systématique cette structure fondamentale : \cle{How + adjectif}, qui sert à poser des questions de mesure en anglais.

\courssub{Observer avant de comprendre}

\begin{textebox}[At the health centre in Bafoussam]
Nurse: Good morning. Before the doctor sees you, I need some information. How old are you?

Patient: I am seventeen years old.

Nurse: And how tall are you?

Patient: I am one metre sixty-five tall.

Nurse: How heavy are you? I mean, what is your weight?

Patient: About sixty kilograms.

Nurse: How long have you had this cough?

Patient: For two weeks now.

Nurse: I see. And how far is your village from here?

Patient: It is about ten kilometres away. I came by bendskin.

Nurse: How often do you come to the health centre?

Patient: Not very often. Maybe twice a year.

Nurse: Well, the doctor will see you soon. By the way, how deep is the well in your village? The doctor suspects the water may be contaminated.

Patient: I don't know exactly, but it is quite deep.
\end{textebox}

\textbf{Glossaire}

\begin{center}
\begin{tabularx}{\linewidth}{|X|X|X|}
\hline
\rowcolor{popblueL} Mot anglais & Nature & Traduction \\
\hline
cough & nom & toux \\
\hline
weight & nom & poids \\
\hline
bendskin & nom (Cameroun) & moto-taxi \\
\hline
well & nom & puits \\
\hline
contaminated & adjectif & contaminé \\
\hline
quite & adverbe & assez \\
\hline
\end{tabularx}
\end{center}

\textbf{Questions d'observation}

\begin{enumerate}
\item Quel mot revient au début de chaque question de l'infirmière ?
\item Où se place l'adjectif (\eng{old}, \eng{tall}, \eng{long}, \eng{far}) par rapport à ce mot ?
\item Quel verbe utilise-t-on dans la réponse \eng{I am seventeen years old} ? Le français utilise-t-il le même verbe ?
\end{enumerate}

\textbf{Réponses attendues :} le mot \eng{how} ouvre chaque question ; l'adjectif se colle immédiatement après \eng{how} ; l'anglais utilise \eng{to be} là où le français utilise « avoir » (\emph{j'ai dix-sept ans} → \eng{I am seventeen}).

\courssub{La règle : forme de la question de mesure}

\begin{propriete}[Structure de « How + adjectif »]
\cle{How + adjectif} sert à interroger sur une \textbf{mesure} ou un \textbf{degré} (âge, taille, distance, durée, poids, etc.). L'adjectif se place \textbf{immédiatement après How}, avant l'auxiliaire :

\begin{center}
\textbf{How + adjectif + auxiliaire + sujet + verbe ?}
\end{center}

\eng{How tall is Mount Cameroon?} (Quelle est la hauteur du mont Cameroun ?)

\eng{How long does the journey take?} (Combien de temps dure le trajet ?)
\end{propriete}

Deux cas selon le verbe :

\begin{itemize}
\item Avec \cle{to be} : pas d'auxiliaire, on inverse simplement sujet et verbe. \eng{How old \textbf{is she}?} (et jamais \eng{How old does she be?}).
\item Avec les autres verbes : on utilise l'auxiliaire \cle{do/does/did}. \eng{How long \textbf{does} the treatment \textbf{last}?} (Combien de temps dure le traitement ?)
\end{itemize}

\begin{popfigure}
\begin{center}
\begin{tikzpicture}[
box/.style={draw=popblue, fill=popblueL, rounded corners=3pt, align=center, minimum height=0.9cm, font=\small},
lbl/.style={font=\footnotesize\itshape, text=popmuted},
arr/.style={-{Stealth[length=2.5mm]}, thick, poporange}]
\node[box] (how) {\textbf{How}};
\node[box, right=0.45cm of how, align=center] (adj){\textbf{adjectif}\\ old, tall, far...};
\node[box, right=0.45cm of adj, align=center] (aux){\textbf{auxiliaire}\\ is / does / did};
\node[box, right=0.45cm of aux] (subj) {\textbf{sujet}};
\node[box, right=0.45cm of subj] (verb) {\textbf{verbe} ?};
\draw[arr] (how) -- (adj);
\draw[arr] (adj) -- (aux);
\draw[arr] (aux) -- (subj);
\draw[arr] (subj) -- (verb);
\node[lbl, below=0.35cm of how] (e1) {How};
\node[lbl, below=0.35cm of adj] (e2) {far};
\node[lbl, below=0.35cm of aux] (e3) {is};
\node[lbl, below=0.35cm of subj] (e4) {the hospital};
\node[lbl, below=0.35cm of verb] (e5) {?};
\draw[poporange, dashed] (how.south) -- (e1.north);
\draw[poporange, dashed] (adj.south) -- (e2.north);
\draw[poporange, dashed] (aux.south) -- (e3.north);
\draw[poporange, dashed] (subj.south) -- (e4.north);
\draw[poporange, dashed] (verb.south) -- (e5.north);
\end{tikzpicture}
\end{center}
\legende{Schéma de la structure « How + adjectif » avec l'exemple annoté \eng{How far is the hospital?}}
\end{popfigure}

\courssub{Les combinaisons essentielles à connaître}

Chaque dimension possède son adjectif attitré. Il faut les mémoriser comme des blocs.

\begin{center}
\begin{tabularx}{\linewidth}{|X|X|X|}
\hline
\rowcolor{popblueL} Combinaison & Mesure demandée & Exemple traduit \\
\hline
How old & âge & \eng{How old is your father?} Quel âge a ton père ? \\
\hline
How tall & taille (personnes, arbres, bâtiments) & \eng{How tall are you?} Combien mesures-tu ? \\
\hline
How high & hauteur (montagnes, murs, objets) & \eng{How high is the wall?} Quelle est la hauteur du mur ? \\
\hline
How long & durée ou longueur & \eng{How long is the Wouri bridge?} Quelle est la longueur du pont sur le Wouri ? \\
\hline
How far & distance & \eng{How far is Douala from Yaoundé?} À quelle distance est Douala de Yaoundé ? \\
\hline
How big & taille, superficie & \eng{How big is your farm?} Quelle est la taille de ta ferme ? \\
\hline
How deep & profondeur & \eng{How deep is the well?} Quelle est la profondeur du puits ? \\
\hline
How wide & largeur & \eng{How wide is the river?} Quelle est la largeur de la rivière ? \\
\hline
How heavy & poids & \eng{How heavy is this bag of cocoa?} Combien pèse ce sac de cacao ? \\
\hline
How often & fréquence & \eng{How often do you exercise?} À quelle fréquence fais-tu du sport ? \\
\hline
\end{tabularx}
\end{center}

\begin{aretenir}[Distinctions délicates]
\begin{itemize}
\item \cle{How tall} pour les personnes (et ce qui « se dresse » : arbres, immeubles) ; \cle{How high} pour les montagnes et les objets mesurés du sol au sommet. \eng{How tall is your brother?} mais \eng{How high is Mount Cameroon?}
\item \cle{How long} = durée OU longueur ; \cle{How far} = distance entre deux lieux. On ne dit jamais \eng{How far does the film last?}
\item \cle{How often} n'est pas un adjectif de mesure mais un adverbe de fréquence ; il suit la même construction.
\end{itemize}
\end{aretenir}

\courssub{Les réponses : to be + mesure}

\begin{propriete}[Construire la réponse]
On répond avec \cle{to be + mesure} :

\eng{She is twenty years old.} (Elle a vingt ans.)

\eng{He is 1.75 m tall.} (Il mesure 1,75 m.)

\eng{It is 5 km away.} (C'est à 5 km.)

Remarques :
\begin{itemize}
\item Pour l'âge, \eng{years old} est facultatif : \eng{She is twenty.} est correct.
\item Pour la distance, on ajoute souvent \cle{away} : \eng{The hospital is three kilometres away.}
\item Pour la durée, on répond avec \cle{for + durée} : \eng{For two weeks.}
\end{itemize}
\end{propriete}

\begin{exemplebox}[Questions et réponses]
\eng{How old is your grandmother?} (Quel âge a ta grand-mère ?) — \eng{She is seventy years old.} (Elle a soixante-dix ans.)

\eng{How far is the hospital?} (À quelle distance est l'hôpital ?) — \eng{It is three kilometres away.} (Il est à trois kilomètres.)

\eng{How long does the treatment last?} (Combien de temps dure le traitement ?) — \eng{Two weeks.} (Deux semaines.)

\eng{How deep is this river?} (Quelle est la profondeur de cette rivière ?) — \eng{It is about four metres deep.} (Elle fait environ quatre mètres de profondeur.)

\eng{How often does the bus pass?} (À quelle fréquence passe le bus ?) — \eng{Every thirty minutes.} (Toutes les trente minutes.)
\end{exemplebox}

\courssub{Comparaison avec le français}

\begin{center}
\begin{tabularx}{\linewidth}{|X|X|X|}
\hline
\rowcolor{popblueL} Français & English & Remarque \\
\hline
Quel âge a-t-elle ? & How old is she? & Le français dit « avoir » un âge ; l'anglais « être » un âge. \\
\hline
Combien mesure-t-il ? & How tall is he? & Pas de verbe « mesurer » obligatoire en anglais : \eng{to be} suffit. \\
\hline
Combien de temps dure le film ? & How long does the film last? & Auxiliaire \eng{does} car le verbe est \eng{last}. \\
\hline
À quelle distance est le marché ? & How far is the market? & « À quelle distance » se condense en \eng{how far}. \\
\hline
Combien pèse ce carton ? & How heavy is this box? & On peut aussi dire \eng{How much does this box weigh?} \\
\hline
\end{tabularx}
\end{center}

\begin{attention}[Erreurs fréquentes des francophones]
\begin{itemize}
\item \eng{She has 20 years.} FAUX (calque de « elle a 20 ans »). En anglais : \eng{She is 20 (years old).} L'âge s'exprime avec \cle{to be}, jamais avec \cle{to have}.
\item \eng{How old have you?} FAUX. On dit \eng{How old are you?}
\item \eng{How long is the film last?} FAUX (double verbe). On dit \eng{How long does the film last?} ou \eng{How long is the film?}
\item Ne pas confondre \eng{How long?} (durée/longueur) et \eng{How far?} (distance) : \eng{How far is Buea from Limbe?} demande une distance, pas une durée.
\item Attention à l'ordre des mots : \eng{How tall is your brother?} et non \eng{How tall your brother is?} (l'inversion est obligatoire dans la question directe).
\end{itemize}
\end{attention}

\courssub{« How + adjectif » dans les exclamations}

La même structure sert aussi à exprimer l'admiration ou la surprise. Dans ce cas, \textbf{pas d'inversion} : le sujet reste avant le verbe, et la phrase se termine par un point d'exclamation.

\begin{propriete}[Exclamation avec How]
\textbf{How + adjectif + sujet + verbe !}

\eng{How beautiful this landscape is!} (Comme ce paysage est beau !)

\eng{How high this mountain is!} (Que cette montagne est haute !)

Compare : question \eng{How high is this mountain?} (inversion) / exclamation \eng{How high this mountain is!} (ordre normal).
\end{propriete}

\begin{exoresolu}[Transformations]
\textbf{Énoncé.} Transforme chaque phrase en question avec \eng{How + adjectif}.

1. My uncle is forty-five years old. (Ask about his age.)

2. The Sanaga River is very wide. (Ask about its width.)

3. The journey from Yaoundé to Bamenda takes six hours. (Ask about the duration.)

4. The stadium is eight kilometres from here. (Ask about the distance.)

\tcblower
\textbf{Solution détaillée.}

1. \eng{How old is your uncle?} — L'âge se demande avec \eng{how old} ; verbe \eng{to be}, donc simple inversion, sans auxiliaire.

2. \eng{How wide is the Sanaga River?} — La largeur se demande avec \eng{how wide} ; inversion de \eng{is}.

3. \eng{How long does the journey from Yaoundé to Bamenda take?} — Le verbe est \eng{take} (pas \eng{to be}), donc auxiliaire \eng{does} + base verbale \eng{take}.

4. \eng{How far is the stadium from here?} — La distance entre deux lieux se demande avec \eng{how far} (et non \eng{how long}, qui concerne la durée ou la longueur).
\end{exoresolu}

\begin{exercice}[À toi de jouer]
1. Pose la question correspondant à chaque réponse :

a) \eng{She is sixteen years old.}

b) \eng{It is two kilometres away.}

c) \eng{He is 1.80 m tall.}

d) \eng{The well is twelve metres deep.}

2. Traduis en anglais :

a) Quel âge a ta sœur ?

b) À quelle distance est l'école ?

c) Combien de temps dure la consultation ?

d) Comme ce jardin est beau !
\end{exercice}

\begin{corrige}[Exercice « À toi de jouer »]
1. a) \eng{How old is she?} b) \eng{How far is it?} c) \eng{How tall is he?} d) \eng{How deep is the well?}

2. a) \eng{How old is your sister?} b) \eng{How far is the school?} c) \eng{How long does the consultation last?} d) \eng{How beautiful this garden is!}
\end{corrige}

\begin{savaistu}[Feet and inches]
Au Royaume-Uni et aux États-Unis, la taille se dit en \cle{feet and inches} (pieds et pouces) : \eng{He is six feet tall} signifie qu'il mesure environ 1,83 m (un foot = environ 30 cm ; un inch = 2.54cm). On prononce « six feet tall » : « siks fite tol ». Au Bac camerounais, la réponse en mètres (\eng{He is 1.75 m tall}) est parfaitement acceptée.
\end{savaistu}

Nous savons désormais interroger sur une mesure avec \eng{How + adjectif}. Dans la section suivante, nous verrons comment comparer et situer un élément au sommet d'une échelle grâce au \textbf{superlatif}.

\coursec{Le superlatif : formation et emplois}

Dans la section précédente, nous avons appris à poser des questions de mesure avec \eng{How + adjective} : \eng{How tall is Mount Cameroon?} Nous savons donc demander la taille, la profondeur ou la chaleur d'un élément. Mais souvent, nous ne voulons pas seulement mesurer : nous voulons \cle{classer}, désigner l'élément qui arrive \emph{en tête} du groupe. Quel est le mois le plus sec ? Quel est l'hôpital le mieux équipé du pays ? Pour répondre à ces questions, l'anglais utilise une structure précise : le \cle{superlatif}. C'est l'objet de cette section.

\courssub{Observation : le superlatif en contexte}

Lisons d'abord ce court article de presse sur la santé et l'environnement au Cameroun. Repérez les adjectifs précédés de \eng{the}.

\begin{textebox}[Health and Environment Report — Yaoundé]
Cameroon is a country of extremes. Mount Cameroon, near Buea, is the highest peak in West Africa, and it is also one of the most active volcanoes on the continent. In the Far North region, April is usually the hottest month of the year, while January is the driest. Health workers say that malaria remains the deadliest disease in the country, especially among young children. “It is the worst health problem we have ever faced,” explains Dr. Ngono, who works in one of the busiest hospitals in Douala. Fortunately, the situation is improving: last year was the best year for vaccination campaigns, and the newest clinics are now the cleanest and the most modern in the region. However, the most serious challenge is access to clean water in the most remote villages. “The furthest communities are often the poorest,” adds Dr. Ngono. “They need the greatest attention.”
\end{textebox}

\textbf{Glossaire}

\begin{center}
\begin{tabularx}{\linewidth}{|l|l|X|}
\hline
\rowcolor{popblueL} Mot anglais & Nature & Traduction \\
\hline
peak & nom & sommet, pic \\
\hline
deadly & adjectif & mortel, meurtrier \\
\hline
to face & verbe & affronter \\
\hline
remote & adjectif & isolé, reculé \\
\hline
access & nom & accès \\
\hline
\end{tabularx}
\end{center}

\textbf{Questions de compréhension}

\begin{enumerate}
\item What is the highest peak in West Africa?
\item According to Dr. Ngono, what is the worst health problem they have ever faced?
\item Which communities need the greatest attention, and why?
\end{enumerate}

Relevons maintenant les formes superlatives du texte : \eng{the highest}, \eng{the most active}, \eng{the hottest}, \eng{the driest}, \eng{the deadliest}, \eng{the worst}, \eng{the busiest}, \eng{the best}, \eng{the newest}, \eng{the cleanest}, \eng{the most modern}, \eng{the most serious}, \eng{the most remote}, \eng{the furthest}, \eng{the greatest}. On remarque deux grandes familles : les adjectifs qui prennent \cle{-est} et ceux qui prennent \cle{the most}. Voyons la règle qui gouverne ce choix.

\courssub{La règle générale}

\begin{propriete}[Le superlatif : définition et principe]
Le superlatif exprime le degré \cle{le plus élevé} d'une qualité au sein d'un groupe : il désigne l'élément qui dépasse tous les autres. Il est \textbf{toujours précédé de} \eng{the} :
\begin{itemize}
\item \eng{the most serious problem} « le problème le plus grave »
\item \eng{the tallest building} « le bâtiment le plus haut »
\end{itemize}
L'anglais choisit entre deux marqueurs selon la \textbf{longueur de l'adjectif} :
\begin{itemize}
\item adjectifs \textbf{courts} (une syllabe) : \eng{the + adjectif + -est} ;
\item adjectifs \textbf{longs} (deux syllabes et plus) : \eng{the most + adjectif}.
\end{itemize}
\end{propriete}

\begin{attention}[Un seul marqueur, et toujours « the »]
Ne jamais dire \eng{the more beautiful} ni \eng{the most biggest} : un seul marqueur de superlatif à la fois ! Et n'oubliez jamais l'article : \eng{Mount Cameroon is highest peak} est faux ; il faut \eng{Mount Cameroon is the highest peak}. De même, ne traduisez pas « le plus » par \eng{the more} : « le plus grand » se dit \eng{the biggest} ou \eng{the tallest}, jamais \eng{the more big}.
\end{attention}

\courssub{Les adjectifs courts : the + adjectif + -est}

\begin{propriete}[Adjectifs d'une syllabe]
Pour les adjectifs d'une syllabe, on ajoute \cle{-est} à l'adjectif, précédé de \eng{the} :
\eng{clean → the cleanest} (le plus propre), \eng{tall → the tallest} (le plus grand), \eng{high → the highest} (le plus haut), \eng{new → the newest} (le plus récent).
\end{propriete}

Trois cas d'orthographe exigent une attention particulière :

\begin{propriete}[Modifications orthographiques]
\begin{enumerate}
\item \textbf{Consonne finale après voyelle courte : on double la consonne.} \eng{big → the biggest} (le plus gros), \eng{hot → the hottest} (le plus chaud), \eng{thin → the thinnest} (le plus mince).
\item \textbf{Adjectifs en -y : le y devient i + -est.} \eng{healthy → the healthiest} (le plus sain), \eng{dirty → the dirtiest} (le plus sale), \eng{busy → the busiest} (le plus chargé), \eng{dry → the driest} (le plus sec).
\item \textbf{Adjectifs en -e muet : on ajoute simplement -st.} \eng{large → the largest} (le plus grand), \eng{safe → the safest} (le plus sûr).
\end{enumerate}
\end{propriete}

\begin{exemplebox}[Exemples traduits]
\eng{January is the driest month in the North.} « Janvier est le mois le plus sec dans le Nord. »

\eng{The Sanaga is the longest river in Cameroon.} « La Sanaga est le plus long fleuve du Cameroun. »

\eng{This is the cleanest hospital in the region.} « C'est l'hôpital le plus propre de la région. »
\end{exemplebox}

\courssub{Les adjectifs longs : the most + adjectif}

\begin{propriete}[Adjectifs de deux syllabes et plus]
Pour les adjectifs longs, on place \cle{the most} devant l'adjectif, qui reste \textbf{invariable} :
\eng{the most polluted} (le plus pollué), \eng{the most beautiful} (le plus beau), \eng{the most dangerous} (le plus dangereux), \eng{the most serious} (le plus grave), \eng{the most expensive} (le plus cher).
\end{propriete}

\begin{exemplebox}[Exemples traduits]
\eng{Douala is one of the most polluted cities in the region.} « Douala est l'une des villes les plus polluées de la région. »

\eng{Cholera is the most dangerous disease in this area.} « Le choléra est la maladie la plus dangereuse de cette zone. »
\end{exemplebox}

\textbf{Cas particulier des adjectifs de deux syllabes.} Certains adjectifs de deux syllabes acceptent les deux formes, surtout ceux en \eng{-y}, \eng{-er}, \eng{-ow}, \eng{-le} : \eng{happy → the happiest}, \eng{narrow → the narrowest}, \eng{simple → the simplest}. En revanche, les participes passés employés comme adjectifs prennent toujours \eng{the most} : \eng{the most polluted}, \eng{the most crowded}, jamais \eng{the pollutedest}.

\courssub{Les superlatifs irréguliers}

\begin{propriete}[Trois irréguliers à connaître par cœur]
\begin{itemize}
\item \eng{good → the best} (bon → le meilleur)
\item \eng{bad → the worst} (mauvais → le pire)
\item \eng{far → the farthest / the furthest} (loin → le plus loin)
\end{itemize}
\end{propriete}

\begin{savaistu}[Farthest ou furthest ?]
\eng{The farthest} et \eng{the furthest} sont tous deux corrects. On préfère \eng{farthest} pour les distances physiques (\eng{the farthest village from the main road}) et \eng{furthest} au sens figuré (\eng{the furthest point of the discussion}). Bonne nouvelle : le français partage les irrégularités \eng{good/best} et \eng{bad/worst} — « bon / le meilleur », « mauvais / le pire » !
\end{savaistu}

\courssub{Tableau récapitulatif de formation}

\begin{popfigure}
\begin{tikzpicture}
\node[draw=popblue, fill=popblueL, rounded corners, font=\bfseries, minimum width=15.5cm, minimum height=0.8cm] at (0,0) {Formation du superlatif anglais};
\node[draw=popline, fill=popcream, rounded corners, align=left, anchor=north west, text width=15cm] at (-7.75,-0.6)
{\textbf{Adjectifs courts (1 syllabe)} : the + adj. + -est \quad clean $\rightarrow$ the \textbf{cleanest}\\
\textbf{Voyelle courte + consonne} : on double la consonne \quad big $\rightarrow$ the \textbf{biggest}\\
\textbf{Adjectifs en -y} : y $\rightarrow$ i + -est \quad healthy $\rightarrow$ the \textbf{healthiest}\\
\textbf{Adjectifs longs (2 syll. et +)} : the most + adj. \quad polluted $\rightarrow$ the \textbf{most polluted}\\
\textbf{Irréguliers} : good $\rightarrow$ the \textbf{best} ; bad $\rightarrow$ the \textbf{worst} ; far $\rightarrow$ the \textbf{farthest/furthest}};
\end{tikzpicture}
\legende{Les quatre familles de formation du superlatif, avec un exemple par ligne.}
\end{popfigure}

\begin{center}
\begin{tabularx}{\linewidth}{|X|X|X|}
\hline
\rowcolor{popblueL} Français & English & Remarque \\
\hline
le plus propre & the cleanest & adjectif court : -est \\
\hline
le plus chaud & the hottest & on double le t \\
\hline
le plus sain & the healthiest & y devient i \\
\hline
le plus pollué & the most polluted & adjectif long : the most \\
\hline
le meilleur & the best & irrégulier (good) \\
\hline
le pire & the worst & irrégulier (bad) \\
\hline
le plus loin & the farthest / furthest & irrégulier (far) \\
\hline
\end{tabularx}
\end{center}

\courssub{Emplois et structures utiles}

\begin{propriete}[Comparer un élément à tout un groupe]
Le superlatif sert à situer un élément au sommet d'un ensemble. Il est suivi de :
\begin{itemize}
\item \cle{in + lieu ou groupe} : \eng{the best hospital in Cameroon} « le meilleur hôpital du Cameroun » ; \eng{the tallest student in the class} « l'élève le plus grand de la classe » ;
\item \cle{of + période ou groupe} : \eng{the hottest month of the year} « le mois le plus chaud de l'année » ; \eng{the youngest of the three brothers} « le plus jeune des trois frères ».
\end{itemize}
\textbf{Attention} : après un superlatif, on dit \eng{in} pour les lieux, pas \eng{of} : \eng{the best school in Yaoundé}, jamais \eng{of Yaoundé}.
\end{propriete}

\begin{propriete}[One of the + superlatif + nom pluriel]
Pour dire « l'un des plus... », la structure est : \eng{one of the + superlatif + nom au PLURIEL} :
\eng{Malaria is one of the deadliest diseases in Africa.} « Le paludisme est l'une des maladies les plus meurtrières d'Afrique. »
Ne dites jamais \eng{one of the deadliest disease} : le nom doit être au pluriel, car l'élément appartient à un ensemble.
\end{propriete}

\begin{propriete}[Le superlatif avec le present perfect]
Après un superlatif, l'expérience vécue s'exprime au \cle{present perfect}, souvent avec \eng{ever} :
\eng{It is the worst flood I have ever seen.} « C'est la pire inondation que j'aie jamais vue. »
\eng{This is the best clinic I have ever visited.} « C'est la meilleure clinique que j'aie jamais visitée. »
Le français utilise ici le subjonctif (« que j'aie jamais vue »), l'anglais emploie simplement le present perfect.
\end{propriete}

\begin{exoresolu}[Transformez au superlatif]
Énoncé : complétez chaque phrase avec le superlatif de l'adjectif entre parenthèses.
\begin{enumerate}
\item Mount Cameroon is \ldots\ (high) mountain in West Africa.
\item This is \ldots\ (bad) road accident I have ever witnessed.
\item Yaoundé is one of \ldots\ (green) capitals in Central Africa.
\item April is \ldots\ (hot) month of the year in Maroua.
\end{enumerate}
\tcblower
Solution détaillée :
\begin{enumerate}
\item \eng{the highest} : adjectif d'une syllabe, on ajoute -est, sans oublier \eng{the}. « Le mont Cameroun est la plus haute montagne d'Afrique de l'Ouest. »
\item \eng{the worst} : superlatif irrégulier de \eng{bad}. Le present perfect \eng{have ever witnessed} confirme la structure d'expérience vécue. « C'est le pire accident de la route que j'aie jamais vu. »
\item \eng{the greenest} : structure \eng{one of the + superlatif + nom pluriel} (\eng{capitals} est bien au pluriel). « Yaoundé est l'une des capitales les plus vertes d'Afrique centrale. »
\item \eng{the hottest} : voyelle courte + consonne, on double le t. « Avril est le mois le plus chaud de l'année à Maroua. »
\end{enumerate}
\end{exoresolu}

\begin{attention}[Pièges des francophones]
\begin{itemize}
\item Ne doublez pas les marqueurs : \eng{the most biggest} est impossible ; dites \eng{the biggest}.
\item N'oubliez jamais \eng{the} : \eng{She is best student} est faux → \eng{She is the best student}.
\item Après \eng{one of the + superlatif}, le nom est au pluriel : \eng{one of the most dangerous diseases}.
\item Pour les lieux, utilisez \eng{in}, pas \eng{of} : \eng{the cleanest city in Cameroon}.
\item Ne confondez pas comparatif et superlatif : \eng{better} (meilleur, entre deux) / \eng{the best} (le meilleur, dans un groupe).
\end{itemize}
\end{attention}

\begin{exercice}[À vous de jouer]
Traduisez en anglais :
\begin{enumerate}
\item Le paludisme est l'une des maladies les plus dangereuses du monde.
\item C'est la pire sécheresse que nous ayons jamais connue.
\item Limbe possède les plages les plus propres du pays.
\end{enumerate}
\end{exercice}

\begin{corrige}[Exercice « À vous de jouer »]
\begin{enumerate}
\item \eng{Malaria is one of the most dangerous diseases in the world.} (adjectif long : \eng{the most dangerous} ; nom pluriel \eng{diseases} après \eng{one of the}.)
\item \eng{It is the worst drought we have ever experienced.} (irrégulier \eng{bad → the worst} ; present perfect après le superlatif.)
\item \eng{Limbe has the cleanest beaches in the country.} (adjectif court : \eng{clean → the cleanest} ; \eng{in} pour le lieu.)
\end{enumerate}
\end{corrige}

\begin{aretenir}[L'essentiel du superlatif]
Le superlatif désigne le degré le plus élevé d'une qualité dans un groupe. Il prend toujours \eng{the}. Adjectifs courts : \eng{the + -est} (avec doublement de consonne et y → i si nécessaire) ; adjectifs longs : \eng{the most + adjectif}. Trois irréguliers : \eng{good → the best}, \eng{bad → the worst}, \eng{far → the farthest/furthest}. Structures clés : \eng{in + lieu}, \eng{of + période}, \eng{one of the + superlatif + nom pluriel}, et le present perfect avec \eng{ever} pour l'expérience vécue.
\end{aretenir}

\coursec{Too...to et enough...to : exprimer la limite}

Dans la section précédente, nous avons appris à situer un élément au sommet d'une échelle grâce au superlatif (\eng{the most polluted river in the region}). Mais dans la vie quotidienne — et surtout dans les sujets d'environnement et de santé du chapitre 3 — nous avons souvent besoin d'exprimer autre chose : une \cle{limite}. L'eau est-elle assez propre pour être bue ? La pollution est-elle trop forte pour qu'on puisse respirer normalement ? L'anglais possède deux structures très précises pour cela : \eng{too...to} et \eng{enough...to}. Ce sont des structures \textbf{incontournables au Baccalauréat}, tant en compréhension de texte qu'en expression écrite.

\courssub{Observation : découvrir les structures en contexte}

\begin{textebox}[A health report from Douala]
Last month, health workers visited several neighbourhoods of Douala to test the quality of drinking water. Their conclusions were worrying. In some areas, the water is \textbf{too dirty to drink}, and families have to buy bottled water, which is \textbf{too expensive for many parents to afford}. “The situation is serious,” explains Dr Mbarga. “The pipes are \textbf{too old to carry clean water} to every house. We need new equipment, but the budget is not \textbf{large enough to cover} all the repairs.”

In the schools of the city, the problem is similar. Classrooms are \textbf{too crowded for teachers to control}, and some children are \textbf{too weak to walk} long distances to school because they do not eat \textbf{enough food}. However, there is hope. A local association says the population is \textbf{aware enough of the problem to act}. “Our volunteers are \textbf{old enough to work} and motivated enough to succeed,” says their leader, Amina, 19. “The river near our quarter was once \textbf{too polluted to swim in}; today, after two years of cleaning, the water is \textbf{clean enough to use} for washing. If everyone helps, Douala will soon be healthy enough for all its children.”

\emph{Glossaire :} worrying (adj.) : inquiétant ; pipes (n.) : tuyaux, canalisations ; to afford (v.) : avoir les moyens d'acheter ; crowded (adj.) : bondé ; aware (adj.) : conscient ; quarter (n.) : quartier.
\end{textebox}

\textbf{Questions d'observation :}
\begin{enumerate}
\item Relevez toutes les structures avec \eng{too} et avec \eng{enough}. Où se place \eng{enough} par rapport à l'adjectif ?
\item Dans la phrase \eng{The water is too dirty to drink}, peut-on boire l'eau ? La structure a-t-elle un sens positif ou négatif ?
\item Dans la phrase \eng{The water is clean enough to use}, peut-on utiliser l'eau ?
\end{enumerate}

\courssub{Structure 1 : too + adjectif + to-infinitive}

\begin{propriete}[Too + adjectif + to-infinitive : l'excès qui empêche]
\textbf{Forme :} \cle{too} + adjectif (+ \eng{for} + complément) + \eng{to} + base verbale.

\textbf{Emploi :} cette structure exprime un \textbf{excès} qui rend l'action \textbf{impossible ou déconseillée}. Elle a un \textbf{sens implicite négatif} : ce qui suit \eng{to} ne peut PAS se faire.

\eng{The box is too heavy to carry.} « La caisse est trop lourde pour être portée » → \textbf{on ne peut pas la porter}.

\textbf{Équivalent français :} « trop... pour (+ infinitif) ».
\end{propriete}

\begin{exemplebox}[Exemples traduits]
\eng{The river is too polluted to swim in.} « La rivière est trop polluée pour s'y baigner. »

\eng{The child is too weak to walk.} « L'enfant est trop faible pour marcher. »

\eng{This exercise is too difficult for the students to do.} « Cet exercice est trop difficile pour que les élèves le fassent. »

\eng{It is too late to catch the last bus to Bamenda.} « Il est trop tard pour attraper le dernier bus pour Bamenda. »
\end{exemplebox}

\begin{attention}[Ne confondez pas too et very !]
\eng{Too} signifie « trop » et implique toujours un \textbf{problème}. \eng{Very} signifie « très » et reste \textbf{neutre}.

\eng{The soup is very hot.} « La soupe est très chaude. » (simple constat : on peut la manger en soufflant dessus)

\eng{The soup is too hot to eat.} « La soupe est trop chaude pour être mangée. » (conséquence : impossible de la manger maintenant)

Piège classique : dire \eng{The food is too good} pour faire un compliment est une erreur ! Cela signifie que la nourriture est « trop bonne », donc qu'il y a un problème. Pour un compliment, dites : \eng{The food is very good.} ou \eng{The food is really good.}
\end{attention}

\courssub{Structure 2 : adjectif + enough + to-infinitive}

\begin{propriete}[Adjectif + enough + to-infinitive : la suffisance qui permet]
\textbf{Forme :} adjectif + \cle{enough} (+ \eng{for} + complément) + \eng{to} + base verbale.

\textbf{Emploi :} cette structure exprime une \textbf{qualité suffisante} qui rend l'action \textbf{possible}. Elle a un sens \textbf{positif}.

\eng{She is old enough to vote.} « Elle est assez âgée pour voter. » → elle PEUT voter.

\textbf{Règle de place capitale :} \eng{enough} se place \textbf{APRÈS l'adjectif} (\eng{tall enough}) mais \textbf{AVANT le nom} (\eng{enough water}, \eng{enough time}). Il se place aussi APRÈS un adverbe : \eng{He speaks English well enough to work in Buea.} « Il parle anglais assez bien pour travailler à Buea. »
\end{propriete}

\begin{exemplebox}[Exemples traduits]
\eng{This water is clean enough to drink.} « Cette eau est assez propre pour être bue. »

\eng{He is strong enough to carry the bag.} « Il est assez fort pour porter le sac. »

\eng{The room is big enough for thirty students.} « La salle est assez grande pour trente élèves. »

\eng{We don't have enough time to finish the project.} « Nous n'avons pas assez de temps pour finir le projet. » (ici \eng{enough} précède le nom \eng{time})
\end{exemplebox}

\begin{attention}[Le piège majeur : la place de enough]
Les francophones placent souvent \eng{enough} AVANT l'adjectif, par calque du français « assez grand ». C'est \textbf{faux} en anglais !

\begin{center}
\begin{tabularx}{\linewidth}{|X|X|}\hline
\rowcolor{popblueL} \textbf{Faux} & \textbf{Correct} \\ \hline
He is enough tall to play basketball. & He is \textbf{tall enough} to play basketball. \\ \hline
The water is enough clean. & The water is \textbf{clean enough}. \\ \hline
\end{tabularx}
\end{center}

En revanche, avec un \textbf{nom}, \eng{enough} se place AVANT : \eng{enough money} (assez d'argent), \eng{enough food} (assez de nourriture), \eng{enough sleep} (assez de sommeil).
\end{attention}

\courssub{Tableau récapitulatif et schéma de transformation}

\begin{center}
\begin{tabularx}{\linewidth}{|X|X|X|}\hline
\rowcolor{popblueL} \textbf{Structure} & \textbf{Sens} & \textbf{Exemple} \\ \hline
too + adj + to V & excès → impossible (négatif) & The air is too polluted to breathe. \\ \hline
adj + enough + to V & suffisance → possible (positif) & The air is clean enough to breathe. \\ \hline
not + adj + enough + to V & insuffisance → impossible (négatif) & He is not old enough to drive. \\ \hline
too + adj + for + n/pron + to V & précise qui est concerné & The text is too hard for me to read. \\ \hline
too much + nom indénombrable & trop de... & too much pollution, too much noise \\ \hline
too many + nom dénombrable & trop de... & too many cars, too many diseases \\ \hline
\end{tabularx}
\end{center}

\begin{popfigure}
\begin{tikzpicture}[node distance=0.6cm]
\node[draw=popink, fill=poppinkL, rounded corners, align=center, text width=5.2cm, minimum height=1.6cm] (too) {\textbf{too + ADJ + to V}\\ sens négatif\\ \eng{He is too young to drive.}};
\node[draw=popgreen, fill=popgreenL, rounded corners, align=center, text width=5.2cm, minimum height=1.6cm, right=1.4cm of too] (enough) {\textbf{ADJ + enough + to V}\\ sens positif\\ \eng{He is old enough to drive.}};
\node[draw=poporange, fill=poporangeL, rounded corners, align=center, text width=5.2cm, below=1.4cm of too] (neg) {\textbf{not + ADJ contraire + enough + to V}\\ \eng{He is not old enough to drive.}};
\draw[-{Stealth}, very thick, poporange] (too) -- (neg) node[midway, left, align=center, text width=2.4cm] {\footnotesize adjectif contraire + négation};
\draw[-{Stealth}, very thick, popgreen] (neg.east) -- (enough.south west);
\draw[{Stealth}-{Stealth}, thick, popblue] (too) -- (enough) node[midway, above] {\footnotesize transformation};
\end{tikzpicture}
\legende{Les deux structures et la transformation too $\leftrightarrow$ enough avec adjectif contraire.}
\end{popfigure}

\begin{methode}[Transformer too...to en not...enough to]
Pour transformer une phrase sans changer son sens :
\begin{enumerate}
\item \textbf{Repérez l'adjectif} : \eng{The shirt is too \textbf{small} for me.}
\item \textbf{Prenez son contraire} : small $\leftrightarrow$ big ; young $\leftrightarrow$ old ; weak $\leftrightarrow$ strong.
\item \textbf{Passez de \eng{too...to} à \eng{not...enough to}} : \eng{The shirt is \textbf{not big enough} for me.}
\end{enumerate}
Vérification : \eng{He is too young to drive} = \eng{He is not old enough to drive.} (Il est trop jeune pour conduire = il n'est pas assez âgé pour conduire.)
\end{methode}

\begin{savaistu}[Préciser qui est concerné : too...for + nom + to]
La structure \eng{too + adjectif + for + nom/pronom + to-infinitive} permet de dire \textbf{qui} subit la limite, ce que le français rend souvent par « pour que » + subjonctif : \eng{The exercise is too difficult for the students to do.} « L'exercice est trop difficile pour que les élèves le fassent. » De même : \eng{The bag is light enough for the child to carry.} « Le sac est assez léger pour que l'enfant le porte. » Attention : après \eng{for}, on utilise le pronom \textbf{complément} : \eng{for me, for him, for them} — jamais \eng{for I} ou \eng{for they}. Cette tournure est très appréciée à l'écrit au Baccalauréat !
\end{savaistu}

\courssub{Entraînement}

\begin{exoresolu}[Transformation guidée]
Énoncé : transformez la phrase suivante avec \eng{enough}, sans changer le sens : \eng{The water is too dirty to drink.}
\tcblower
Solution détaillée :
\begin{enumerate}
\item Adjectif repéré : \eng{dirty}. Son contraire : \eng{clean}.
\item On applique la structure \eng{not + adjectif contraire + enough + to} : la saleté empêche de boire, donc l'eau « n'est pas assez propre pour être bue ».
\item Résultat : \eng{The water is \textbf{not clean enough} to drink.}
\end{enumerate}
Vérification du sens : dans les deux phrases, on ne peut pas boire l'eau. La transformation est correcte.
\end{exoresolu}

\begin{exercice}[À vous de jouer]
\begin{enumerate}
\item Complétez avec \eng{too} ou \eng{enough} : a) This coffee is \dots{} hot to drink. b) Are you warm \dots{}? c) We don't have \dots{} money to buy the medicine. d) He is \dots{} young to travel alone.
\item Transformez avec \eng{enough} : a) The child is too weak to walk. b) The room is too small for forty students.
\item Traduisez : « Il y a trop de voitures et trop de pollution à Douala. »
\item Traduisez : « Cet hôpital est assez grand pour accueillir tous les malades. »
\end{enumerate}
\end{exercice}

\begin{corrige}[Exercice « À vous de jouer »]
1. a) \eng{too} (excès : on ne peut pas le boire) ; b) \eng{enough} (après l'adjectif \eng{warm}) ; c) \eng{enough} (avant le nom \eng{money}) ; d) \eng{too} (excès qui empêche).

2. a) \eng{The child is not strong enough to walk.} b) \eng{The room is not big enough for forty students.}

3. \eng{There are too many cars and too much pollution in Douala.} (\eng{cars} est dénombrable → \eng{too many} ; \eng{pollution} est indénombrable → \eng{too much}.)

4. \eng{This hospital is big enough to receive all the patients.} (\eng{enough} après l'adjectif \eng{big} ; \eng{patients} = malades, faux ami de « patients » au sens de « personnes patientes ».)
\end{corrige}

\begin{aretenir}[L'essentiel de la section]
\begin{itemize}
\item \eng{too + adjectif + to V} = excès qui \textbf{empêche} (sens négatif) : \eng{too dirty to drink}.
\item \eng{adjectif + enough + to V} = qualité qui \textbf{permet} (sens positif) : \eng{clean enough to drink}.
\item \eng{enough} APRÈS l'adjectif ou l'adverbe, AVANT le nom : \eng{old enough}, \eng{fast enough}, mais \eng{enough time}.
\item Transformation : \eng{too young to drive} = \eng{not old enough to drive}.
\item \eng{too} = problème ; \eng{very} = simple intensité. \eng{too much} + indénombrable, \eng{too many} + dénombrable.
\end{itemize}
\end{aretenir}

\coursec{Observation : découvrir les structures en contexte}

\begin{textebox}[Dialogue : Une visite au centre de santé de Buea]
\eng{— Good morning, Nurse Agnes. How \textbf{old} is this health centre?} \\
\eng{— Good morning. It is \textbf{twenty years old}. It was built in 2005.} \\
\eng{— And \textbf{how far} is the nearest hospital?} \\
\eng{— The referral hospital in Buea town is \textbf{only five kilometres away}, but the road is \textbf{too bad} for cars during the rainy season. Patients \textbf{too weak} to walk must wait for a motorbike.} \\
\eng{— Is the water \textbf{clean enough} to drink?} \\
\eng{— Not really. The well is \textbf{not deep enough}. The water is \textbf{too dirty}. Actually, this centre is \textbf{the most crowded} in the district, but it has \textbf{the fewest} nurses.} \\
\eng{— \textbf{How deep} is the well?} \\
\eng{— About \textbf{ten metres}. The \textbf{deepest} well in the village is fifteen metres, but it dried up last year.}
\end{textebox}

\begin{exemplebox}[Vocabulaire utile du dialogue (glossaire)]
\begin{itemize}
\item \eng{referral hospital} (n.) : hôpital de référence
\item \eng{rainy season} (n.) : saison des pluies
\item \eng{motorbike} (n.) : moto-taxi (bendskin)
\item \eng{well} (n.) : puits
\item \eng{crowded} (adj.) : bondé, surpeuplé
\item \eng{dried up} (v., past) : s'est asséché
\item \eng{district} (n.) : district, arrondissement
\end{itemize}
\end{exemplebox}

\begin{methode}[Comment observer un texte pour en extraire la grammaire]
\begin{enumerate}
\item \textbf{Lis / écoute} le texte une première fois pour le sens global (qui parle ? de quoi ? où ?).
\item \textbf{Surligne} toutes les structures qui comparent, mesurent ou expriment une limite : \eng{How + adjectif}, \eng{the + superlatif}, \eng{too + adjectif + to}, \eng{adjectif + enough + to}.
\item \textbf{Classe-les} dans un tableau à trois colonnes : \textit{Question de mesure} / \textit{Superlatif} / \textit{Limite (too / enough)}.
\item \textbf{Déduis} la règle : forme (ordre des mots) et emploi (sens).
\end{enumerate}
\end{methode}

\coursec{How + adjectif : les questions de mesure}

\begin{propriete}[Structure et emploi de \eng{How + adjectif}]
\eng{How} s'associe à un adjectif pour demander une \textbf{mesure précise} (âge, distance, dimension, quantité). La structure est invariable :
\begin{center}
\eng{How + adjectif + auxiliaire + sujet + verbe ?}
\end{center}
\begin{itemize}
\item L'adjectif suit immédiatement \eng{How} (pas de nom entre les deux).
\item La réponse se fait par une \textbf{valeur chiffrée + unité} (souvent omise si le contexte est clair).
\end{itemize}
\end{propriete}

\begin{center}
\begin{tabularx}{\linewidth}{|X|X|X|}
\hline
\rowcolor{popblueL} \textbf{Question type} & \textbf{Exemple} & \textbf{Réponse type} \\
\hline
\eng{How old} (âge) & \eng{How old is your brother?} & \eng{He is 17 (years old).} \\
\hline
\eng{How far} (distance) & \eng{How far is Douala from Yaoundé?} & \eng{It is 250 km.} \\
\hline
\eng{How deep} (profondeur) & \eng{How deep is the well?} & \eng{It is 10 metres deep.} \\
\hline
\eng{How long} (longueur / durée) & \eng{How long is the bridge? / How long is the film?} & \eng{500 metres. / Two hours.} \\
\hline
\eng{How high / tall} (hauteur) & \eng{How high is Mount Cameroon?} & \eng{4,095 metres.} \\
\hline
\eng{How much / many} (quantité) & \eng{How much water? / How many nurses?} & \eng{20 litres. / Three nurses.} \\
\hline
\end{tabularx}
\end{center}

\begin{attention}[Piège : \eng{How old} vs \eng{How many years}]
\begin{itemize}
\item ✓ \eng{How old are you?} (standard)
\item ✗ \eng{How many years do you have?} (calque du français « Quel âge as-tu ? »)
\item ✗ \eng{How many years old are you?} (redondant)
\end{itemize}
\eng{Old} est l'adjectif qui porte la notion d'âge ; \eng{years} n'apparaît que dans la réponse.
\end{attention}

\begin{exercice}[Questions \eng{How + adjectif} — Mise en pratique]
À partir des réponses, écris la question exacte :
\begin{enumerate}
\item \ldots\ ? — My father is fifty-two years old.
\item \ldots\ ? — The river is three metres deep here.
\item \ldots\ ? — The nearest pharmacy is two kilometres away.
\item \ldots\ ? — The Harmattan season lasts three months.
\item \ldots\ ? — Mount Cameroon is 4,095 metres high.
\end{enumerate}
\end{exercice}

\begin{corrige}[Exercice \eng{How + adjectif}]
1. \eng{How old is your father?} ; 2. \eng{How deep is the river (here)?} ; 3. \eng{How far is the nearest pharmacy?} ; 4. \eng{How long does the Harmattan season last?} (ou \eng{How long is the Harmattan season?}) ; 5. \eng{How high is Mount Cameroon?}
\end{corrige}

\coursec{Le superlatif : formation et emplois}

\begin{propriete}[Règle de formation du superlatif de supériorité]
Le superlatif compare \textbf{un élément à tout un groupe} (le plus / le moins). Il se forme de deux façons selon la longueur de l'adjectif :
\begin{enumerate}
\item \textbf{Adjectifs courts (1 syllabe, ou 2 syllabes en \eng{-y}, \eng{-er}, \eng{-ow}, \eng{-le})} : \eng{the + adjectif + -est}.
\item \textbf{Adjectifs longs (2 syllabes et plus, sauf exceptions)} : \eng{the most + adjectif}.
\item \textbf{Adjectifs irréguliers} : \eng{good $\to$ the best} ; \eng{bad $\to$ the worst} ; \eng{far $\to$ the farthest / the furthest} ; \eng{little $\to$ the least} ; \eng{much/many $\to$ the most}.
\end{enumerate}
\textbf{Toujours} placer l'article défini \eng{the} devant le superlatif.
\end{propriete}

\begin{center}
\begin{tabular}{|l|l|l|l|}
\hline
\rowcolor{popblueL} \textbf{Adjectif} & \textbf{Type} & \textbf{Règle} & \textbf{Superlatif} \\
\hline
\eng{clean} & court & +\eng{-est} & \eng{the cleanest} \\
\hline
\eng{hot} & court (CVC) & double cons. +\eng{-est} & \eng{the hottest} \\
\hline
\eng{dirty} & 2 syll. en \eng{-y} & \eng{y $\to$ i} +\eng{-est} & \eng{the dirtiest} \\
\hline
\eng{narrow} & 2 syll. en \eng{-ow} & +\eng{-est} & \eng{the narrowest} \\
\hline
\eng{simple} & 2 syll. en \eng{-le} & +\eng{-est} & \eng{the simplest} \\
\hline
\eng{polluted} & long (3 syll.) & \eng{the most} + adj. & \eng{the most polluted} \\
\hline
\eng{dangerous} & long & \eng{the most} + adj. & \eng{the most dangerous} \\
\hline
\eng{good / bad} & irrégulier & -- & \eng{the best / the worst} \\
\hline
\end{tabular}
\end{center}

\begin{attention}[Erreur fatale : le double marqueur]
\begin{itemize}
\item ✗ \eng{the most cleanest}, \eng{the most dirtiest}, \eng{the most hottest}.
\item ✓ \eng{the cleanest}, \eng{the dirtiest}, \eng{the hottest}.
\end{itemize}
Un seul marqueur de superlatif : \textbf{soit} \eng{-est} \textbf{soit} \eng{most}, jamais les deux.
\end{attention}

\begin{exemplebox}[Superlatif dans le contexte camerounais]
\begin{itemize}
\item \eng{Yaoundé is \textbf{the most populated} city in Cameroon.} (adjectif long)
\item \eng{Mount Cameroon is \textbf{the highest} peak in West Africa.} (adjectif court \eng{high} $\to$ \eng{highest})
\item \eng{The dry season is \textbf{the hottest} period in Maroua.} (CVC \eng{hot} $\to$ \eng{hottest})
\item \eng{Dr. Nkeng is \textbf{the best} surgeon in the region.} (irrégulier)
\item \eng{This clinic has \textbf{the least} equipment of all.} (irrégulier \eng{little $\to$ least})
\end{itemize}
\end{exemplebox}

\begin{exercice}[Formation du superlatif]
Mets l'adjectif entre parenthèses au superlatif de supériorité :
\begin{enumerate}
\item The Wouri bridge is \ldots\ (long) bridge in Central Africa.
\item Malaria is \ldots\ (serious) health problem in our region.
\item My grandmother is \ldots\ (old) person in the village.
\item This medicine is \ldots\ (effective) of all.
\item The road to Bamenda is \ldots\ (bad) during the rains.
\end{enumerate}
\end{exercice}

\begin{corrige}[Exercice Superlatif]
1. \eng{the longest} (court) ; 2. \eng{the most serious} (long) ; 3. \eng{the oldest} (court) ; 4. \eng{the most effective} (long) ; 5. \eng{the worst} (irrégulier \eng{bad}).
\end{corrige}

\coursec{Too...to et enough...to : exprimer la limite}

\begin{propriete}[\eng{Too + adjectif + to} : excès négatif (impossibilité)]
\eng{Too} signifie « trop » (dépassement d'une norme, conséquence négative). Structure :
\eng{Sujet + be + too + adjectif + to + infinitif}
\eng{The patient is too weak to walk.} « Le patient est trop faible pour marcher » $\Rightarrow$ Il \textbf{ne peut pas} marcher.
\end{propriete}

\begin{propriete}[\eng{Adjectif + enough + to} : suffisance (capacité)]
\eng{Enough} signifie « assez » (atteinte de la norme, conséquence positive). Structure :
\eng{Sujet + be + adjectif + enough + to + infinitif}
\eng{The well is deep enough to supply the village.} « Le puits est assez profond pour approvisionner le village » $\Rightarrow$ Il \textbf{peut} le faire.
\end{propriete}

\begin{propriete}[\eng{Not + adjectif + enough + to} : insuffisance (impossibilité)]
La forme négative de \eng{enough} exprime la même chose que \eng{too} + contraire.
\eng{The water is not clean enough to drink.} = \eng{The water is too dirty to drink.}
\end{propriete}

\begin{center}
\begin{tabularx}{\linewidth}{|X|X|X|}
\hline
\rowcolor{popblueL} \textbf{Structure} & \textbf{Sens} & \textbf{Exemple (Contexte santé/environnement)} \\
\hline
\eng{too + adj + to} & Excès $\to$ Impossible & \eng{The smoke is too thick to breathe.} \\
\hline
\eng{adj + enough + to} & Suffisant $\to$ Possible & \eng{The vaccine is stable enough to transport.} \\
\hline
\eng{not + adj + enough + to} & Insuffisant $\to$ Impossible & \eng{The clinic is not big enough to host everyone.} \\
\hline
\eng{enough + nom + to} & Quantité suffisante & \eng{We have enough beds to admit them.} \\
\hline
\end{tabularx}
\end{center}

\begin{attention}[La place de \eng{enough} : le piège n°1 des francophones]
\begin{itemize}
\item \textbf{Devant un adjectif / adverbe} : \eng{enough} se place \textbf{APRÈS}.
  \eng{rich enough}, \eng{quickly enough}, \eng{clean enough}.
  (Français : « assez riche » $\to$ \textbf{avant}).
\item \textbf{Devant un nom} : \eng{enough} se place \textbf{AVANT}.
  \eng{enough water}, \eng{enough time}, \eng{enough nurses}.
  (Français : « assez d'eau » $\to$ \textbf{avant}).
\end{itemize}
\end{attention}

\begin{exemplebox}[Paires \eng{too} / \eng{enough} (transformations)]
\begin{itemize}
\item \eng{The boy is too short to reach the shelf.} $\leftrightarrow$ \eng{The boy is not tall enough to reach the shelf.}
\item \eng{The medicine is too expensive for us.} $\leftrightarrow$ \eng{The medicine is not cheap enough for us.}
\item \eng{The river is not clean enough to swim in.} $\leftrightarrow$ \eng{The river is too dirty to swim in.}
\end{itemize}
\end{exemplebox}

\begin{exercice}[Choix multiple : \eng{too}, \eng{very}, \eng{enough}]
Choisis le mot juste :
\begin{enumerate}
\item The sun in Maroua is \ldots\ hot in April; people stay indoors. (\eng{too} / \eng{very})
\item The nurse was \ldots\ kind to explain the treatment twice. (\eng{too} / \eng{very})
\item Is the water cold \ldots\ to keep the fish fresh? (\eng{enough} / \eng{too})
\item We have \ldots\ mosquito nets for the whole family. (\eng{enough} / \eng{too})
\item The pollution in Douala is \ldots\ dangerous for children. (\eng{too} / \eng{very})
\end{enumerate}
\end{exercice}

\begin{corrige}[Exercice Choix multiple]
1. \eng{too} (problème : on reste enfermé) ; 2. \eng{very} (constat positif, pas de conséquence négative) ; 3. \eng{enough} (après adjectif, sens capacité) ; 4. \eng{enough} (devant nom) ; 5. \eng{too} (danger = problème / excès négatif).
\end{corrige}

\coursec{Pièges des francophones et entraînement ciblé}

Nous venons de voir comment \eng{too...to} et \eng{enough...to} permettent d'exprimer la limite et la capacité. Mais ces structures, comme le superlatif et les questions en \eng{How + adjectif}, sont de véritables pièges pour les élèves francophones : le français nous pousse à traduire mot à mot, et c'est là que naissent les erreurs. Cette dernière partie de la leçon est consacrée à un bilan complet des erreurs types, à leur correction, puis à un entraînement ciblé pour que ces structures deviennent des réflexes — notamment au moment de l'épreuve écrite du Baccalauréat.

\courssub{Bilan des erreurs types des francophones}

\begin{attention}[Les quatre erreurs emblématiques — à bannir définitivement]
Voici les quatre fautes les plus fréquentes dans les copies camerounaises. Chacune vient d'une traduction littérale du français.
\begin{itemize}
\item \textbf{L'âge avec \eng{to have}.} ✗ \eng{My sister has 15 years.} → ✓ \eng{My sister is 15 (years old).} En anglais, l'âge se dit avec \eng{to be}, jamais avec \eng{to have}. Le français dit « ma sœur a 15 ans », l'anglais dit littéralement « ma sœur est 15 ans ».
\item \textbf{Le double superlatif.} ✗ \eng{the most tall} → ✓ \eng{the tallest}. Un adjectif court prend \eng{-est} ; un adjectif long prend \eng{most}. Mais jamais les deux en même temps : un seul marqueur de superlatif.
\item \textbf{La place de \eng{enough}.} ✗ \eng{enough rich} → ✓ \eng{rich enough}. En français, « assez » se place \emph{avant} l'adjectif (« assez riche ») ; en anglais, \eng{enough} se place \emph{après} l'adjectif. Attention : devant un nom, \eng{enough} revient avant : \eng{enough water} (« assez d'eau »).
\item \textbf{La confusion \eng{too} / \eng{very}.} ✗ \eng{too much hot} → ✓ \eng{too hot}. \eng{Too} signifie « trop » (un problème, un excès négatif) ; \eng{very} signifie « très » (une simple intensité). On ne dit jamais \eng{too much} devant un adjectif seul.
\end{itemize}
\end{attention}

\begin{popfigure}
\begin{tikzpicture}[scale=0.9, every node/.style={font=\small}]
% Headers
\node[fill=popblue, text=white, font=\bfseries, minimum width=7.2cm, minimum height=0.8cm, align=center] (h1) at (0,0) {Erreur fréquente (✗)};
\node[fill=popgreen, text=white, font=\bfseries, minimum width=7.2cm, minimum height=0.8cm, align=center] (h2) at (7.6,0) {Correction (✓)};
% Row 1
\node[fill=popblueL, minimum width=7.2cm, minimum height=0.9cm, align=center] at (0,-0.9) {\eng{My sister has 15 years.}};
\node[fill=popgreenL, minimum width=7.2cm, minimum height=0.9cm, align=center] at (7.6,-0.9) {\eng{My sister is 15 (years old).}};
% Row 2
\node[fill=popblueL, minimum width=7.2cm, minimum height=0.9cm, align=center] at (0,-1.9) {\eng{the most tall building}};
\node[fill=popgreenL, minimum width=7.2cm, minimum height=0.9cm, align=center] at (7.6,-1.9) {\eng{the tallest building}};
% Row 3
\node[fill=popblueL, minimum width=7.2cm, minimum height=0.9cm, align=center] at (0,-2.9) {\eng{enough rich to buy it}};
\node[fill=popgreenL, minimum width=7.2cm, minimum height=0.9cm, align=center] at (7.6,-2.9) {\eng{rich enough to buy it}};
% Row 4
\node[fill=popblueL, minimum width=7.2cm, minimum height=0.9cm, align=center] at (0,-3.9) {\eng{The water is too much hot.}};
\node[fill=popgreenL, minimum width=7.2cm, minimum height=0.9cm, align=center] at (7.6,-3.9) {\eng{The water is too hot.}};
% Lines
\draw[popline, thick] (-3.6,-0.45) -- (11.2,-0.45);
\draw[popline] (-3.6,-1.4) -- (11.2,-1.4);
\draw[popline] (-3.6,-2.4) -- (11.2,-2.4);
\draw[popline] (-3.6,-3.4) -- (11.2,-3.4);
\draw[popline] (-3.6,-4.4) -- (11.2,-4.4);
\end{tikzpicture}
\legende{Tableau « Erreur / Correction » : les quatre fautes emblématiques des francophones.}
\end{popfigure}

\courssub{Correction collective : un corpus de phrases d'élèves camerounais}

Voici des phrases réellement écrites par des élèves de Terminale lors de devoirs sur l'environnement et la santé. Pour chacune, identifions la faute, expliquons-la, puis corrigeons.

\begin{exoresolu}[Correction d'une copie d'élève]
Phrase écrite par un élève de Yaoundé : \eng{Yaoundé is the most polluted city because the taxis are too much old and the streets are not enough clean.}

Corrige cette phrase en identifiant toutes les erreurs.
\tcblower
\textbf{Solution détaillée.} La phrase contient trois erreurs (une structure est correcte) :
\begin{enumerate}
\item \eng{the most polluted} est \textbf{correct} : \eng{polluted} est un adjectif long (trois syllabes : pol-lu-ted), donc le superlatif avec \eng{most} est juste. Il ne faut pas « corriger » ce qui est bien !
\item ✗ \eng{too much old} → ✓ \eng{too old}. Devant un adjectif seul, on utilise \eng{too}, jamais \eng{too much}. (\eng{Too much} ne s'emploie que devant un nom indénombrable : \eng{too much noise} « trop de bruit », \eng{too much pollution}.)
\item ✗ \eng{not enough clean} → ✓ \eng{not clean enough}. \eng{Enough} se place \emph{après} l'adjectif.
\end{enumerate}
\textbf{Phrase corrigée :} \eng{Yaoundé is the most polluted city because the taxis are too old and the streets are not clean enough.}

Traduction : « Yaoundé est la ville la plus polluée parce que les taxis sont trop vieux et que les rues ne sont pas assez propres. »
\end{exoresolu}

\begin{exemplebox}[D'autres phrases du corpus, corrigées et traduites]
\begin{itemize}
\item ✗ \eng{My grandmother has 70 years and she is too much weak to walk.} → ✓ \eng{My grandmother is 70 (years old) and she is too weak to walk.} « Ma grand-mère a 70 ans et elle est trop faible pour marcher. »
\item ✗ \eng{Mount Cameroon is the most high mountain in West Africa.} → ✓ \eng{Mount Cameroon is the highest mountain in West Africa.} « Le mont Cameroun est la plus haute montagne d'Afrique de l'Ouest. » (\eng{high} = court $\to$ \eng{-est}).
\item ✗ \eng{The well is not enough deep for the village.} → ✓ \eng{The well is not deep enough for the village.} « Le puits n'est pas assez profond pour le village. »
\item ✗ \eng{How many old is your school?} → ✓ \eng{How old is your school?} « Quel âge a ton école ? » (\eng{How old}, sans \eng{many}).
\item ✗ \eng{This exercise is too easy to do.} (Contexte : l'élève veut dire qu'il \textbf{peut} le faire). → ✓ \eng{This exercise is easy enough to do.} \eng{Too} implique l'échec ; ici c'est une réussite, donc \eng{enough}.
\end{itemize}
\end{exemplebox}

\begin{center}
\begin{tabularx}{\linewidth}{|X|X|X|}
\hline
\rowcolor{popblueL} \textbf{Français (piège)} & \textbf{English (correct)} & \textbf{Remarque} \\
\hline
Ma sœur a 15 ans. & \eng{My sister is 15 (years old).} & Âge avec \eng{to be}, pas \eng{to have}. \\
\hline
La plus haute montagne & \eng{the highest mountain} & Un seul marqueur : \eng{-est}. \\
\hline
Assez riche pour acheter & \eng{rich enough to buy} & \eng{Enough} après l'adjectif. \\
\hline
Assez d'eau potable & \eng{enough drinking water} & \eng{Enough} avant le nom. \\
\hline
L'eau est trop chaude. & \eng{The water is too hot.} & \eng{Too} = problème ; jamais \eng{too much} + adjectif. \\
\hline
La ville est très chaude. & \eng{The city is very hot.} & \eng{Very} = simple intensité, neutre. \\
\hline
Quelle est la profondeur du puits ? & \eng{How deep is the well?} & \eng{How + adjectif}, pas de nom. \\
\hline
\end{tabularx}
\end{center}

\courssub{Exercices d'application}

\begin{exercice}[Complétion à trous]
Complète avec \eng{too}, \eng{very}, \eng{enough}, \eng{how}, \eng{most} ou \eng{-est} (forme fléchie) :
\begin{enumerate}
\item The hospital is \ldots\ far from our quarter; patients walk for two hours.
\item This clinic is \ldots\ clean: the nurses wash it every day. (simple constat positif)
\item Is the vaccine cold \ldots\ to be kept for a week?
\item \ldots\ long is the bridge over the Wouri river?
\item The Harmattan is the dry\ldots\ season of the year in the Far North.
\end{enumerate}
\end{exercice}

\begin{exercice}[Formation du superlatif]
Mets l'adjectif entre parenthèses au superlatif de supériorité :
\begin{enumerate}
\item Douala is \ldots\ (industrial) city in Cameroon.
\item January is \ldots\ (hot) month in Maroua.
\item Dr. Amina is \ldots\ (good) nurse in the district hospital.
\item Air pollution is one of \ldots\ (serious) problems in big cities.
\end{enumerate}
\end{exercice}

\begin{exercice}[Questions en \eng{How + adjectif}]
Forme la question qui correspond à la réponse :
\begin{enumerate}
\item \ldots\ ? — The well is twenty metres deep.
\item \ldots\ ? — My grandfather is eighty years old.
\item \ldots\ ? — The new market is three kilometres from here.
\end{enumerate}
\end{exercice}

\begin{corrige}[Exercices d'application]
\textbf{Complétion à trous :} 1. \eng{too} (un problème : deux heures de marche) ; 2. \eng{very} (simple intensité positive) ; 3. \eng{enough} (après l'adjectif \eng{cold}) ; 4. \eng{How} ; 5. \eng{driest} (adjectif court en \eng{-y} $\to$ \eng{-iest}).

\textbf{Superlatif :} 1. \eng{the most industrial} ; 2. \eng{the hottest} (doublement du \eng{t}) ; 3. \eng{the best} (irrégulier) ; 4. \eng{the most serious}.

\textbf{Questions :} 1. \eng{How deep is the well?} ; 2. \eng{How old is your grandfather?} ; 3. \eng{How far is the new market from here?}
\end{corrige}

\courssub{Exercices de transformation}

\begin{methode}[Transformer \eng{too} en \eng{enough} (et inversement)]
Pour transformer une phrase sans changer son sens :
\begin{enumerate}
\item Repère l'adjectif de la phrase de départ (ex. \eng{young}).
\item Remplace-le par son \textbf{contraire} (\eng{young} $\leftrightarrow$ \eng{old}, \eng{hot} $\leftrightarrow$ \eng{cold}, \eng{difficult} $\leftrightarrow$ \eng{easy}, \eng{dirty} $\leftrightarrow$ \eng{clean}, \eng{ill} $\leftrightarrow$ \eng{well}).
\item Change la structure : \eng{too + adjectif + to} devient \eng{not + adjectif + enough + to}.
\item Vérifie la place de \eng{enough} : toujours \emph{après} l'adjectif.
\end{enumerate}
Exemple : \eng{He is too young to drive.} $\to$ \eng{He is not old enough to drive.} « Il est trop jeune pour conduire » = « il n'est pas assez âgé pour conduire ».
\end{methode}

\begin{exoresolu}[Transformation guidée]
Transforme : \eng{The water is too dirty to drink.}
\tcblower
\textbf{Étape 1 :} adjectif = \eng{dirty} (« sale »). \textbf{Étape 2 :} contraire = \eng{clean} (« propre »). \textbf{Étape 3 :} \eng{too dirty to} devient \eng{not clean enough to}. \textbf{Étape 4 :} \eng{enough} est bien après \eng{clean}.

\textbf{Résultat :} \eng{The water is not clean enough to drink.} « L'eau n'est pas assez propre pour être bue. »
\end{exoresolu}

\begin{exercice}[À ton tour : transformations]
\begin{enumerate}
\item \eng{The patient is too ill to travel.} $\to$ (transforme avec \eng{enough} et \eng{well})
\item \eng{This exercise is too difficult for me.} $\to$ (transforme avec \eng{enough} et \eng{easy})
\item \eng{She is old enough to vote.} $\to$ (transforme avec \eng{too} et \eng{young})
\item \eng{Buea is colder than Douala. Douala is hotter than Bamenda.} $\to$ (réécris avec un superlatif : \eng{Buea is the\ldots})
\end{enumerate}
\end{exercice}

\begin{corrige}[Exercice de transformations]
1. \eng{The patient is not well enough to travel.} (« Le patient n'est pas assez bien portant pour voyager. ») ; 2. \eng{This exercise is not easy enough for me.} ; 3. \eng{She is too young to vote.} ; 4. \eng{Buea is the coldest of the three towns.} (« Buea est la plus froide des trois villes. »)
\end{corrige}

\courssub{Production personnelle : la santé et l'environnement de ton quartier}

\begin{experience}[Décris ton quartier avec les structures de la leçon]
Par groupes de trois, préparez cinq phrases sur votre quartier (ou votre village) en utilisant obligatoirement :
\begin{itemize}
\item un \textbf{superlatif} : \eng{The most polluted place in our quarter is\ldots} ;
\item une question en \eng{How + adjectif} : \eng{How far is the nearest health centre?} ;
\item une structure \eng{too\ldots to} : \eng{The stream is too dirty to wash clothes in.} ;
\item une structure \eng{enough\ldots to} : \eng{Our school is not big enough to welcome all the pupils.} ;
\item l'âge avec \eng{to be} : \eng{Our health centre is twenty years old.}
\end{itemize}
Chaque groupe lit ensuite ses phrases à la classe ; les autres groupes vérifient avec la grille ci-dessous et signalent les erreurs. Le groupe qui produit cinq phrases sans faute gagne le titre de « Grammar Champions ».
\end{experience}

\begin{exemplebox}[Modèles de phrases personnelles]
\begin{itemize}
\item \eng{Douala is the most industrial city in Cameroon.} « Douala est la ville la plus industrielle du Cameroun. »
\item \eng{How deep is the well?} « Quelle est la profondeur du puits ? »
\item \eng{The patient is too ill to travel.} « Le patient est trop malade pour voyager. »
\item \eng{The rubbish dump is too close to our houses.} « La décharge est trop proche de nos maisons. »
\item \eng{The water is not clean enough to drink.} « L'eau n'est pas assez propre pour être bue. »
\end{itemize}
\end{exemplebox}

\courssub{Auto-évaluation : la grille de vérification avant de rendre ta copie}

\begin{methode}[Les cinq questions à te poser avant de rendre ta copie]
Avant de rendre toute production écrite (devoir, composition, épreuve du Bac), relis-toi en te posant ces cinq questions :
\begin{enumerate}
\item Y a-t-il \eng{the} devant chaque superlatif ? (✓ \eng{the tallest}, ✗ \eng{tallest} seul)
\item N'y a-t-il qu'\textbf{un seul} marqueur de superlatif : \eng{-est} \emph{ou} \eng{most}, jamais les deux ? (✗ \eng{the most tall})
\item \eng{Enough} est-il placé \textbf{après} l'adjectif et \textbf{avant} le nom ? (✓ \eng{rich enough}, ✓ \eng{enough money})
\item L'âge est-il exprimé avec \eng{to be} ? (✓ \eng{She is 15}, ✗ \eng{She has 15 years})
\item \eng{Too} exprime-t-il bien un \textbf{problème} ? Si tu veux simplement dire « très », utilise \eng{very}. Et jamais \eng{too much} devant un adjectif seul.
\end{enumerate}
Si tu réponds « oui » aux cinq questions, ta copie est propre sur ces structures.
\end{methode}

\begin{aretenir}[L'essentiel de la leçon]
\begin{itemize}
\item \eng{How + adjectif} pour les questions de mesure : \eng{How old / deep / far / long / high\ldots?}
\item Superlatif : \eng{the + -est} (adjectifs courts) ou \eng{the most + adjectif} (adjectifs longs) ; irréguliers : \eng{the best}, \eng{the worst}, \eng{the farthest}.
\item \eng{too + adjectif + to} = « trop\ldots pour » (problème) ; \eng{adjectif + enough + to} = « assez\ldots pour » (capacité).
\item L'âge se dit avec \eng{to be} ; \eng{enough} suit l'adjectif ; un seul marqueur de superlatif ; \eng{too} $\neq$ \eng{very}.
\end{itemize}
\end{aretenir}

\begin{savaistu}[Le superlatif, un atout au Bac]
À l'épreuve d'expression écrite du Baccalauréat, les correcteurs évaluent notamment la « richesse grammaticale » de ta copie. Une description qui ne contient que des phrases simples au présent obtient une note moyenne, même sans faute. En revanche, utiliser au moins \textbf{un superlatif} (\eng{the most serious problem in my town is\ldots}) et \textbf{une structure \eng{too} ou \eng{enough}} (\eng{young people are too poor to see a doctor}) montre que tu maîtrises des structures variées — et cela valorise immédiatement ta production. Entraîne-toi à « glisser » ces deux structures dans chaque composition sur l'environnement ou la santé !
\end{savaistu}

\newpage

\coursec{Activité d'intégration}

\courssub{Objectifs de l'activité}

À la fin de cette activité d'intégration, l'élève sera capable de :

\begin{itemize}
\item mobiliser dans une production authentique les trois structures étudiées dans la leçon : \cle{How + adjectif} pour les questions de mesure, le \cle{superlatif} pour hiérarchiser des faits, et \cle{too...to} / \cle{enough...to} pour exprimer la limite ;
\item comprendre et exploiter des informations factuelles sur la santé et l'environnement de sa localité (compréhension de l'écrit et de l'oral) ;
\item poser des questions d'interview correctement formées et y répondre à l'oral (expression orale en interaction) ;
\item rédiger un court reportage informatif de 8 à 10 lignes, clair et bien organisé (expression écrite) ;
\item travailler en équipe, répartir les rôles et évaluer la production des autres groupes selon des critères précis.
\end{itemize}

\courssub{Situation}

Votre établissement participe à la semaine « Health and Environment Week » organisée par le conseil d'arrondissement. La radio communautaire de votre district, \eng{Mile 6 Community Radio}, a lancé un appel aux lycées : chaque établissement peut envoyer un reportage de deux minutes sur la situation sanitaire et environnementale de la ville. Le meilleur reportage sera diffusé le samedi matin, dans l'émission \eng{Youth Voice}. Votre professeur d'anglais a décidé que votre classe de Terminale A relèverait le défi. Par groupes de quatre, vous allez préparer, présenter puis rédiger votre reportage.

Voici la fiche d'information que la radio a envoyée aux écoles participantes :

\begin{textebox}[Mile 6 Community Radio — Information sheet for schools]
Dear students,

Mile 6 Community Radio invites your school to take part in our programme “Youth Voice”. Your task: prepare a two-minute report on “The health and environment situation in our town”.

Your report must answer questions like: How far is the nearest health centre from the poorest quarter? How long do patients wait before they see a nurse? How clean is the water people drink? How serious is the waste problem in the market?

Tell us what the biggest problem is, which quarter is the cleanest, and what the most urgent action should be. Do not only complain: show what is too difficult today, and what is already good enough to give hope. For example: “The district hospital is too small to receive all the patients” or “The new well is deep enough to provide safe water all year round”.

Speak clearly, give facts, and finish with one strong recommendation. The best report will be broadcast on Saturday at 9 a.m.

Good luck!

The “Youth Voice” team
\end{textebox}

\begin{definition}[Glossaire utile]
\begin{itemize}
\item \eng{health centre} (n.) : centre de santé, CSI
\item \eng{quarter} (n.) : quartier (d'une ville)
\item \eng{waste} (n.) : déchets
\item \eng{well} (n.) : puits
\item \eng{to broadcast} (v.) : diffuser (à la radio) ; irrégulier : \eng{broadcast -- broadcast -- broadcast}
\item \eng{urgent} (adj.) : urgent
\end{itemize}
\end{definition}

\courssub{Consignes}

\textbf{Tâche 1 — Compréhension du document (10 min).} Lisez la fiche de la radio en groupe. Répondez par écrit, en anglais et en phrases complètes : (a) What is the topic of the report? (b) How long must the report be? (c) What kind of questions does the radio suggest? (d) What must the report finish with?

\textbf{Tâche 2 — Collecte des faits (10 min).} Dans votre groupe, notez cinq faits réels (ou réalistes) sur votre ville ou village : distances, durées d'attente, état des marchés, de l'eau, des ordures. Classez-les : quel est \emph{the most serious problem}? Quel est \emph{the cleanest place}? Utilisez le superlatif pour justifier votre choix.

\textbf{Tâche 3 — Préparation des questions d'interview (10 min).} Rédigez au moins deux questions en \cle{How + adjectif} qu'un journaliste poserait à un habitant ou à un infirmier : \eng{How far is the nearest health centre from your house?} ; \eng{How long do patients usually wait?} Attention à l'inversion du sujet et du verbe.

\textbf{Tâche 4 — Rédaction du script oral (15 min).} Écrivez le script de votre reportage de deux minutes. Il doit obligatoirement contenir : deux questions en \eng{How + adjectif}, deux superlatifs différents (\eng{the cleanest quarter}, \eng{the most serious problem}), deux structures \eng{too...to} ou \eng{enough...to}, et une recommandation finale. Répartissez les rôles : présentateur, journaliste, deux témoins interrogés.

\textbf{Tâche 5 — Présentation orale (15 min).} Chaque groupe présente son reportage devant la classe, comme à la radio. Les autres groupes écoutent et remplissent la grille d'évaluation : les structures grammaticales sont-elles correctes ? Les faits sont-ils clairs ? La recommandation est-elle convaincante ?

\textbf{Tâche 6 — Production écrite finale (15 min).} Rédigez votre reportage en 8 à 10 lignes sur une affiche, avec un titre en anglais. Les affiches sont exposées au mur de la salle. La classe vote pour le reportage le plus informatif.

\courssub{Ressources à mobiliser}

\begin{aretenir}[Les trois structures à réutiliser]
\begin{itemize}
\item \cle{How + adjectif} pour les questions de mesure : \eng{How far is the clinic?} « À quelle distance est le dispensaire ? » ; \eng{How long do patients wait?} « Combien de temps les patients attendent-ils ? » ; \eng{How deep is the well?} « Quelle est la profondeur du puits ? » Adjectifs fréquents : \eng{far, long, deep, high, wide, old, big, clean, serious}.
\item Le \cle{superlatif} : adjectif court $\rightarrow$ \eng{the + adj + -est} (\eng{the cleanest, the deepest}) ; adjectif long $\rightarrow$ \eng{the most + adj} (\eng{the most serious problem}) ; irréguliers : \eng{good $\rightarrow$ the best}, \eng{bad $\rightarrow$ the worst}, \eng{far $\rightarrow$ the farthest}.
\item \cle{too + adj + to + base verbale} : « trop... pour » (\eng{The clinic is too small to receive all the patients.}) ; \cle{adj + enough + to + base verbale} : « assez... pour » (\eng{The well is deep enough to provide safe water.}). Attention à la place de \eng{enough} : toujours \emph{après} l'adjectif.
\end{itemize}
\end{aretenir}

\begin{attention}[Pièges à éviter dans votre production]
\begin{itemize}
\item Ne dites jamais \eng{How long is the waiting?} pour « combien de temps attend-on ? » : dites \eng{How long do patients wait?}
\item Pas de double marque : jamais \eng{the most cleanest} ; on choisit \emph{soit} \eng{-est}, \emph{soit} \eng{most}.
\item \eng{Enough} ne se place jamais devant l'adjectif : jamais \eng{enough deep}, mais \eng{deep enough}. (Mais devant un nom : \eng{enough water} « assez d'eau ».)
\item Après \eng{too...to} et \eng{enough...to}, le verbe est à la base verbale, sans \eng{to} supplémentaire ni conjugaison : \eng{too small to receive}, jamais \eng{too small for receiving}.
\item N'oubliez pas l'article \eng{the} devant le superlatif : jamais \eng{cleanest quarter}, mais \eng{the cleanest quarter}.
\end{itemize}
\end{attention}

Lexique thématique utile : \eng{health centre, nurse, patient, disease, malaria, safe water, well, tap, rubbish, waste, market, quarter, pollution, to clean up, to protect, to improve}.

\courssub{Proposition de production et corrigé commenté}

\begin{exoresolu}[Un reportage modèle : « Health and Environment in Nkongsamba »]
Rédigez un reportage de 8 à 10 lignes répondant à l'appel de Mile 6 Community Radio, en utilisant au moins deux questions en \eng{How + adjectif}, deux superlatifs et deux structures \eng{too/enough...to}.
\tcblower
\textbf{Production modèle :}

\eng{Good morning, listeners. This is Youth Voice, reporting from Government High School Nkongsamba. How far is the nearest health centre from the Mbanga Road quarter? It is five kilometres away. And how long do patients wait there? Sometimes three hours. The most serious problem in our town is waste: the central market is the dirtiest place, and the rubbish bins are too small to contain all the waste. However, there is hope. Nkongsamba is the cleanest town in the department this year, and the new well near the stadium is deep enough to provide safe water to two thousand families. Our recommendation: the council should build a larger health centre, because the present one is too small to receive all the patients. Thank you for listening.}

\textbf{Commentaire des choix de langue :}
\begin{itemize}
\item \eng{How far is the nearest health centre...?} et \eng{how long do patients wait there?} : deux questions de mesure en \cle{How + adjectif}, avec inversion correcte du sujet et du verbe (\eng{is}, \eng{do... wait}).
\item \eng{The most serious problem} : superlatif d'un adjectif long avec \eng{the most} ; \eng{the dirtiest place} et \eng{the cleanest town} : superlatifs d'adjectifs courts en \eng{-est} (attention à l'orthographe : \eng{dirty} $\rightarrow$ \eng{dirtiest}, le y devient i devant \eng{-est}).
\item \eng{too small to contain} et \eng{too small to receive} : la limite négative, « trop petites pour contenir », « trop petit pour recevoir » ; \eng{deep enough to provide} : la limite positive, « assez profond pour fournir », avec \eng{enough} bien placé après l'adjectif.
\item Préposition : on dit \eng{the cleanest town in the department} (« dans le département »), et non \emph{of}.
\item Structure du texte : salutation radio, faits chiffrés, problème majeur, note d'espoir, recommandation finale — exactement ce que la fiche de la radio demandait.
\end{itemize}
\end{exoresolu}

\courssub{Bilan de l'activité}

Cette activité a permis de réinvestir l'ensemble des savoir-faire de la leçon dans une situation de communication réaliste et motivante. Vous avez compris un document authentique en anglais (la fiche de la radio), collecté et hiérarchisé des informations grâce au superlatif, formulé des questions d'interview en \eng{How + adjectif}, exprimé des limites et des capacités avec \eng{too...to} et \eng{enough...to}, puis présenté un reportage à l'oral avant de le fixer par écrit. L'évaluation par les pairs et le vote final développent en outre l'esprit critique et la coopération. Pour aller plus loin, vous pouvez enregistrer votre reportage sur téléphone et le comparer au modèle : prononciation claire, débit naturel et exactitude grammaticale sont les trois critères d'un bon journaliste radio. Retenez l'essentiel : \eng{How + adjectif} pour mesurer, \eng{the + superlatif} pour classer, \eng{too/enough...to} pour juger — trois outils simples pour décrire le monde qui vous entoure et proposer des solutions.

\newpage

\coursec{Méthodes et exercices résolus}

\courssub{Observation — découvrir les structures en contexte}

\begin{textebox}[A health report from Buea]
How serious is the water problem in our town? Last month, a team of health workers visited three villages near Buea. They found that the main well in the first village was too deep to repair with simple tools, and the water in the second village was too dirty to drink. “This is the worst situation we have seen in the region,” said the head nurse. “The third village is the cleanest of the three, but it is also the furthest from the health centre, so sick people cannot always reach it in time.” The team concluded that the youngest children were the most exposed to disease. Fortunately, the local council is now rich enough to build two new wells, and the pupils are old enough to help keep them clean.
\end{textebox}

\begin{exemplebox}[Repérage des structures dans le texte]
\begin{itemize}
\item Question de mesure : \eng{How serious is the water problem?} (« Quelle est la gravité du problème de l'eau ? »)
\item Superlatifs : \eng{the worst situation}, \eng{the cleanest}, \eng{the furthest}, \eng{the most exposed}, \eng{the youngest children}.
\item Limite avec \eng{too...to} : \eng{too deep to repair}, \eng{too dirty to drink}.
\item Suffisance avec \eng{enough...to} : \eng{rich enough to build}, \eng{old enough to help}.
\end{itemize}
Observer : \eng{too} précède l'adjectif, \eng{enough} le suit, et chaque superlatif porte l'article \eng{the}.
\end{exemplebox}

\courssub{Fiche méthode 1 — Poser une question de mesure avec \eng{How + adjectif}}

\begin{methode}[Construire une question avec \eng{How + adjectif}]
Pour demander une mesure (taille, âge, longueur, profondeur, distance), l'anglais utilise \cle{How + adjectif + auxiliaire + sujet + verbe ?}. La démarche :
\begin{enumerate}
\item Identifier la dimension demandée : taille (\eng{tall}), âge (\eng{old}), longueur (\eng{long}), profondeur (\eng{deep}), largeur (\eng{wide}), hauteur d'un objet (\eng{high}), distance (\eng{far}), poids (\eng{heavy}).
\item Choisir le bon adjectif : \eng{tall} pour les personnes, \eng{high} pour les montagnes et les bâtiments, \eng{old} pour l'âge.
\item Placer l'auxiliaire avant le sujet : \eng{How tall \textbf{are you}?} et jamais \eng{How tall you are?}.
\item Répondre avec la mesure : \eng{I am 1.75 metres tall.} / \eng{The lake is 20 metres deep.}
\end{enumerate}
Réflexe : le français dit « Quel âge as-tu ? », l'anglais dit \eng{How old are you?} — on ne traduit jamais mot à mot.
\end{methode}

\begin{propriete}[Les paires \eng{How + adjectif} à connaître]
\begin{center}
\begin{tabularx}{\linewidth}{|X|X|X|}\hline
\rowcolor{popblueL} Français & English & Remarque \\ \hline
Quel âge... ? & \eng{How old...?} & jamais \eng{What age} \\ \hline
Quelle taille (personne) ? & \eng{How tall...?} & personnes et arbres \\ \hline
Quelle hauteur (montagne, mur) ? & \eng{How high...?} & choses élevées \\ \hline
Quelle longueur ? & \eng{How long...?} & aussi la durée \\ \hline
Quelle profondeur ? & \eng{How deep...?} & puits, rivières \\ \hline
Quelle largeur ? & \eng{How wide...?} & routes, fleuves \\ \hline
À quelle distance ? & \eng{How far...?} & distances \\ \hline
Quel poids ? & \eng{How heavy...?} & objets, charges \\ \hline
\end{tabularx}
\end{center}
\end{propriete}

\begin{exoresolu}[Questions de mesure — type interrogation écrite]
Énoncé : You are interviewing a health worker in Buea. Write the questions using \eng{How + adjective}.
\begin{enumerate}
\item (âge / the patients / usually / be) ?
\item (profondeur / the well / be) ?
\item (hauteur / Mount Cameroon / be) ?
\item (distance / the hospital / from the village / be) ?
\end{enumerate}
\tcblower
Solution détaillée :
\begin{enumerate}
\item L'âge se demande avec \eng{old}. Auxiliaire \eng{are} avant le sujet : \eng{How old are the patients usually?} (« Quel âge ont généralement les patients ? »).
\item La profondeur d'un puits se demande avec \eng{deep} : \eng{How deep is the well?} (« Quelle est la profondeur du puits ? »).
\item Pour une montagne, on utilise \eng{high} (et non \eng{tall}, réservé aux personnes) : \eng{How high is Mount Cameroon?} (« Quelle est l'altitude du mont Cameroun ? »).
\item La distance se demande avec \eng{far} : \eng{How far is the hospital from the village?} (« À quelle distance l'hôpital se trouve-t-il du village ? »).
\end{enumerate}
Conclusion : retenir les couples \eng{old/âge}, \eng{tall/personnes}, \eng{high/montagnes et bâtiments}, \eng{deep/profondeur}, \eng{far/distance}, et placer toujours l'auxiliaire avant le sujet.
\end{exoresolu}

\courssub{Fiche méthode 2 — Former et employer le superlatif}

\begin{methode}[Le superlatif : \eng{the + adj-est} ou \eng{the most + adj}]
Pour situer un élément au sommet d'un groupe, la démarche est :
\begin{enumerate}
\item Compter les syllabes de l'adjectif.
\item Adjectif court (1 syllabe, parfois 2) : \cle{the + adjectif + -est}. Exemples : \eng{the cleanest}, \eng{the hottest}. Attention aux doubles consonnes : \eng{big} $\rightarrow$ \eng{the biggest} ; \eng{-y} devient \eng{-i} : \eng{healthy} $\rightarrow$ \eng{the healthiest}.
\item Adjectif long (2 syllabes et plus) : \cle{the most + adjectif}. Exemple : \eng{the most polluted city}.
\item Mémoriser les irréguliers : \eng{good} $\rightarrow$ \eng{the best} ; \eng{bad} $\rightarrow$ \eng{the worst} ; \eng{far} $\rightarrow$ \eng{the furthest/the farthest}.
\item Compléter avec \eng{in} (lieu, groupe) ou \eng{of} (période, ensemble) : \eng{the best hospital \textbf{in} Douala}, \eng{the worst flood \textbf{of} the decade}.
\item Ne jamais oublier l'article \cle{the} et ne jamais ajouter \eng{most} à un adjectif déjà en \eng{-est}.
\end{enumerate}
\end{methode}

\begin{propriete}[Tableau récapitulatif des formes du superlatif]
\begin{center}
\begin{tabular}{|l|l|l|}\hline
\rowcolor{popblueL} Type d'adjectif & Formation & Exemple \\ \hline
Court (1 syllabe) & the + adj + -est & \eng{tall} $\rightarrow$ \eng{the tallest} \\ \hline
Court, consonne finale doublée & the + adj + consonne + -est & \eng{big} $\rightarrow$ \eng{the biggest} \\ \hline
Court en -e muet & the + adj + -st & \eng{large} $\rightarrow$ \eng{the largest} \\ \hline
Deux syllabes en -y & the + adj(-y $\rightarrow$ -i) + -est & \eng{dirty} $\rightarrow$ \eng{the dirtiest} \\ \hline
Long (2 syllabes et plus) & the most + adj & \eng{dangerous} $\rightarrow$ \eng{the most dangerous} \\ \hline
Irrégulier & forme spéciale & \eng{good} $\rightarrow$ \eng{the best} ; \eng{bad} $\rightarrow$ \eng{the worst} \\ \hline
\end{tabular}
\end{center}
Exception utile : quelques adjectifs de deux syllabes acceptent les deux formes, par exemple \eng{common} $\rightarrow$ \eng{the commonest} ou \eng{the most common}.
\end{propriete}

\begin{exoresolu}[Superlatifs — type Bac, exercice de transformation]
Énoncé : Complete the sentences with the correct superlative form of the adjective in brackets.
\begin{enumerate}
\item The Mfoundi river is \ldots\ldots\ldots (polluted) river in Yaoundé.
\item Prevention is \ldots\ldots\ldots (good) cure.
\item Last year was \ldots\ldots\ldots (dry) season of the century in the Far North.
\item This is \ldots\ldots\ldots (clean) market in the town.
\end{enumerate}
\tcblower
Solution détaillée :
\begin{enumerate}
\item \eng{Polluted} a trois syllabes $\rightarrow$ \eng{the most polluted} : \eng{The Mfoundi river is the most polluted river in Yaoundé.} (« La rivière Mfoundi est la rivière la plus polluée de Yaoundé. »)
\item \eng{Good} est irrégulier $\rightarrow$ \eng{the best} : \eng{Prevention is the best cure.} (« La prévention est le meilleur remède. ») Ne jamais écrire \eng{the goodest}.
\item \eng{Dry} est court, \eng{-y} devient \eng{-i} $\rightarrow$ \eng{the driest} : \eng{Last year was the driest season of the century in the Far North.}
\item \eng{Clean} est court $\rightarrow$ \eng{the cleanest}, avec l'article \eng{the} : \eng{This is the cleanest market in the town.}
\end{enumerate}
Conclusion : toujours vérifier trois choses — le nombre de syllabes, l'article \eng{the} obligatoire, et les irréguliers \eng{best/worst} appris par cœur.
\end{exoresolu}

\courssub{Fiche méthode 3 — Exprimer la limite avec \eng{too...to} et \eng{enough...to}}

\begin{methode}[Choisir entre \eng{too...to} et \eng{enough...to}]
\begin{enumerate}
\item \cle{too + adjectif + to + verbe} exprime un excès qui empêche l'action (sens négatif) : \eng{The water is too dirty to drink.} (« L'eau est trop sale pour être bue. »)
\item \cle{adjectif + enough + to + verbe} exprime une quantité suffisante qui permet l'action (sens positif) : \eng{The boy is old enough to go to school.} (« Le garçon est assez grand pour aller à l'école. »)
\item Point critique de position : \eng{too} se place \textbf{avant} l'adjectif, \eng{enough} se place \textbf{après} l'adjectif mais \textbf{avant} un nom : \eng{enough food}, \eng{enough clean water}.
\item Pour transformer une phrase : \eng{too + adj} équivaut à \eng{not + adj + enough} avec l'adjectif contraire. \eng{He is too weak to walk.} = \eng{He is not strong enough to walk.}
\item Ajouter un sujet différent avec \eng{for} : \eng{The box is too heavy \textbf{for me} to carry.}
\end{enumerate}
\end{methode}

\begin{aretenir}[Positions de \eng{too} et \eng{enough} en une image]
\begin{itemize}
\item \eng{too} + adjectif : \eng{too hot}, \eng{too expensive}.
\item adjectif + \eng{enough} : \eng{strong enough}, \eng{old enough}.
\item \eng{enough} + nom : \eng{enough money}, \eng{enough time}.
\item \eng{too much} + nom indénombrable : \eng{too much noise} ; jamais devant un adjectif.
\end{itemize}
\end{aretenir}

\begin{exoresolu}[Transformation \eng{too/enough} — type Bac]
Énoncé : Rewrite each sentence using the word in brackets, without changing the meaning.
\begin{enumerate}
\item The air in the city is so polluted that children cannot play outside. (\eng{too})
\item The nurse is very experienced. She can manage the clinic alone. (\eng{enough})
\item The river is too deep for the children to cross. (\eng{enough})
\end{enumerate}
\tcblower
Solution détaillée :
\begin{enumerate}
\item Le sens est « excès qui empêche » $\rightarrow$ \eng{too + adj + to} : \eng{The air in the city is too polluted for children to play outside.} On introduit le sujet \eng{children} avec \eng{for}. (« L'air de la ville est trop pollué pour que les enfants jouent dehors. »)
\item Le sens est « suffisamment... pour » $\rightarrow$ \eng{adj + enough + to} : \eng{The nurse is experienced enough to manage the clinic alone.} Bien placer \eng{enough} après \eng{experienced}.
\item Transformation par le contraire : \eng{deep} $\rightarrow$ contraire \eng{shallow} : \eng{The river is not shallow enough for the children to cross.} La négation \eng{not} est indispensable pour garder le sens.
\end{enumerate}
Conclusion : avant de transformer, demander « excès ou suffisance ? », puis vérifier la position de \eng{too} (avant l'adjectif) et de \eng{enough} (après l'adjectif).
\end{exoresolu}

\courssub{Fiche méthode 4 — Utiliser les structures dans un court texte}

\begin{methode}[Rédiger un paragraphe avec superlatifs, \eng{How + adj} et \eng{too/enough}]
Pour une production écrite type Bac (environnement, santé) :
\begin{enumerate}
\item Ouvrir avec un superlatif pour accrocher : \eng{Malaria is one of the most dangerous diseases in our region.}
\item Varier avec une question de mesure en tête de paragraphe ou en conclusion : \eng{How serious is the problem?}
\item Développer avec \eng{too...to} pour décrire un obstacle : \eng{Some villages are too far from the health centre to reach it on foot.}
\item Nuancer avec \eng{enough...to} pour proposer une solution : \eng{The government has built enough wells to supply clean water.}
\item Relire en vérifiant : article \eng{the} devant chaque superlatif, position de \eng{too/enough}, aucun calque du français.
\end{enumerate}
\end{methode}

\begin{exoresolu}[Production écrite guidée — type Bac]
Énoncé : Write a short paragraph (about 60 words) about sanitation in your town. Use at least one superlative, one \eng{too...to} structure and one \eng{enough...to} structure.
\tcblower
Solution détaillée — démarche et modèle de réponse :
\begin{enumerate}
\item Superlatif d'ouverture : choisir un adjectif long $\rightarrow$ \eng{the most serious}.
\item Obstacle avec \eng{too...to} : \eng{too dirty to let the rain water flow away}.
\item Solution avec \eng{enough...to} : \eng{old enough to help}.
\end{enumerate}
Réponse rédigée : \eng{Poor sanitation is the most serious problem in my town. The gutters are too dirty to let the rain water flow away, so the streets flood every season. How dangerous is this for our health? Very dangerous, because dirty water spreads cholera. Fortunately, the pupils are old enough to help clean the neighbourhood every Saturday morning.} (« Le manque d'assainissement est le problème le plus grave de ma ville. Les caniveaux sont trop sales pour laisser l'eau de pluie s'écouler, donc les rues sont inondées à chaque saison. Quel danger cela représente-t-il pour notre santé ? Un grand danger, car l'eau sale propage le choléra. Heureusement, les élèves sont assez grands pour aider à nettoyer le quartier chaque samedi matin. »)
Conclusion : le correcteur cherche les trois structures exigées, correctement formées (\eng{the most serious}, \eng{too dirty to}, \eng{old enough to}) et intégrées dans un texte cohérent.
\end{exoresolu}

\begin{attention}[Les 5 erreurs qui coûtent des points au Bac]
\begin{enumerate}
\item Oublier l'article devant le superlatif : \eng{He is tallest boy.} $\rightarrow$ \eng{He is \textbf{the} tallest boy.}
\item Doubler la marque du superlatif : \eng{the most biggest} $\rightarrow$ \eng{the biggest}. Jamais \eng{most} avec un adjectif en \eng{-est}.
\item Calquer le français « Quel âge as-tu ? » : \eng{What age are you?} $\rightarrow$ \eng{How old are you?} De même, jamais \eng{How is your age?}.
\item Mal placer \eng{enough} : \eng{enough old to vote} $\rightarrow$ \eng{old \textbf{enough} to vote}. Mais devant un nom : \eng{enough water}.
\item Confondre \eng{too} (excès négatif) et \eng{very} (simple intensité) : \eng{The soup is too hot} signifie qu'on ne peut pas la manger ; pour dire « très chaude » sans idée d'obstacle, écrire \eng{The soup is very hot}. Ne pas traduire « trop » par \eng{too much} devant un adjectif : \eng{too much tired} $\rightarrow$ \eng{too tired}.
\end{enumerate}
\end{attention}

\begin{exercice}[Entraînement ciblé — toutes structures confondues]
\begin{enumerate}
\item Ask questions with \eng{How + adjective} : (a) you want to know the depth of a lake ; (b) you want to know a friend's age ; (c) you want to know the height of a building in Douala.
\item Put the adjectives in brackets into the superlative : (a) This is \ldots\ldots\ldots (bad) road in the region. (b) January is \ldots\ldots\ldots (hot) month of the year here. (c) She is \ldots\ldots\ldots (intelligent) girl in the class.
\item Rewrite with \eng{too} or \eng{enough} : (a) The child is very young. He cannot go to school. (b) The water is so cold that we cannot swim in it. (c) He is strong. He can carry the bag.
\end{enumerate}
\end{exercice}

\begin{corrige}[Entraînement ciblé]
\begin{enumerate}
\item (a) \eng{How deep is the lake?} (b) \eng{How old are you?} / \eng{How old is your friend?} (c) \eng{How high is the building in Douala?}
\item (a) \eng{the worst} (irrégulier) ; (b) \eng{the hottest} (consonne doublée) ; (c) \eng{the most intelligent} (adjectif long).
\item (a) \eng{The child is too young to go to school.} (b) \eng{The water is too cold for us to swim in.} (c) \eng{He is strong enough to carry the bag.}
\end{enumerate}
\end{corrige}

\newpage

\coursec{Exercices}

\courssub{Je vérifie mes connaissances}

\begin{exercice}[QCM : la bonne forme]
Pour chaque phrase, choisis la bonne réponse (A, B ou C) et recopie la phrase complète.
\begin{enumerate}
\item Mount Cameroon is \textbf{...} mountain in the country. \quad A. the higher \quad B. the highest \quad C. the most high
\item \textbf{...} is the Wouri River? — About 160 kilometres. \quad A. How far \quad B. How long \quad C. How much
\item This water is \textbf{...} dirty to drink. \quad A. enough \quad B. too \quad C. very
\item She is old \textbf{...} to vote. \quad A. too \quad B. enough \quad C. so
\item Yaoundé is \textbf{...} than Bamenda in December. \quad A. more cold \quad B. the coldest \quad C. colder
\item \textbf{...} is your grandmother? — She is seventy-five. \quad A. How many years \quad B. How old \quad C. How long
\item That was \textbf{...} interesting lesson of the year. \quad A. the most \quad B. the more \quad C. most
\item The box is too heavy \textbf{...} carry. \quad A. for \quad B. to \quad C. that
\item \textbf{...} is this mango tree? — About ten metres. \quad A. How tall \quad B. How high is tall \quad C. How much tall
\item He speaks English well \textbf{...} to work as a guide. \quad A. too \quad B. enough \quad C. very
\end{enumerate}
\end{exercice}

\begin{exercice}[Vrai ou faux ?]
Indique si chaque phrase est \eng{true} ou \eng{false}, puis justifie ta réponse en citant la règle de grammaire concernée.
\begin{enumerate}
\item \eng{“How old are you?” is the correct way to ask someone's age.}
\item \eng{We say “the most big city” for the superlative of “big”.}
\item \eng{“Too” has a negative meaning: it means “more than is good or possible”.}
\item \eng{In “old enough to drive”, the word “enough” comes before the adjective.}
\item \eng{“How far” is used to ask about distance.}
\item \eng{Short adjectives like “hot” form the superlative with “the most”.}
\end{enumerate}
\end{exercice}

\begin{exercice}[Vocabulaire : les questions de mesure]
Associe chaque question anglaise à ce qu'elle demande. Recopie les paires correctes.
\begin{center}
\begin{tabularx}{\linewidth}{|X|X|}
\hline
\rowcolor{popblueL} \textbf{Question} & \textbf{Ce qu'elle demande} \\
\hline
1. How old is he? & a. the distance between two places \\
\hline
2. How far is Buea from Douala? & b. the age of a person \\
\hline
3. How tall is the building? & c. the price or quantity \\
\hline
4. How long is the bridge? & d. the height of a thing or person \\
\hline
5. How much is this fish? & e. the length of a thing \\
\hline
6. How deep is the lake? & f. the depth of water \\
\hline
\end{tabularx}
\end{center}
\end{exercice}

\begin{exercice}[Superlatifs à compléter]
Complète chaque phrase avec le superlatif de l'adjectif entre parenthèses.
\begin{enumerate}
\item The Wouri estuary is one of \textbf{...} (large) in Central Africa.
\item January is often \textbf{...} (dry) month in the North of Cameroon.
\item This is \textbf{...} (good) hospital in the region.
\item The Dja Reserve is one of \textbf{...} (old) forests in Africa.
\item That was \textbf{...} (bad) flood the town has ever seen.
\item My grandmother is \textbf{...} (wise) person in our family.
\end{enumerate}
\end{exercice}

\courssub{Je m'entraîne}

\begin{exercice}[Traduction]
Traduis les phrases suivantes en anglais, en utilisant les structures de la leçon.
\begin{enumerate}
\item Quelle est la profondeur de ce puits ? — Il fait douze mètres de profondeur.
\item Douala est la ville la plus chaude du pays en mars.
\item Ce sac est trop lourd pour que je le porte.
\item Elle est assez grande pour atteindre l'étagère.
\item Quel âge a ton petit frère ? — Il a dix ans.
\item C'est le marché le plus animé de Yaoundé.
\end{enumerate}
\end{exercice}

\begin{exercice}[Transformation avec \eng{enough}]
Réécris chaque phrase avec \eng{enough} sans changer le sens, comme dans le modèle.\par
\emph{Modèle :} \eng{The room is too dark to read in.} $\rightarrow$ \eng{The room is not bright enough to read in.}
\begin{enumerate}
\item The boy is too young to go to school alone.
\item This soup is too hot to eat now.
\item He is too short to reach the window.
\item The road is too narrow for two lorries to pass.
\item She is too tired to walk to the market.
\end{enumerate}
\end{exercice}

\begin{textebox}[Air quality in Douala]
Douala is one of \textbf{(1)} cities in Central Africa. The air is often \textbf{(2)} dirty for children with asthma to play outside. Near the industrial zone, the smoke is \textbf{(3)} thick for drivers to see clearly. “\textbf{(4)} is the factory from our school?” asked the pupils. “Only two kilometres,” answered the teacher. The level of dust in the air is among \textbf{(5)} ever recorded in the country. Last year was \textbf{(6)} year for respiratory diseases in the city. Many families do not earn \textbf{(7)} money to move to a cleaner neighbourhood. “\textbf{(8)} will this situation last?” wondered a mother. “And \textbf{(9)} does it cost to buy clean water every week?” Her daughter, who is \textbf{(10)} to understand the problem, decided to join the school's environment club.
\end{textebox}

\begin{exercice}[Texte à trous : la pollution à Douala]
Complète le texte ci-dessus avec dix mots ou expressions choisis dans la liste : \eng{the most polluted — how far — too — enough — the highest — how long — the worst — too — how much — old enough}.
\end{exercice}

\begin{exercice}[Construis la question]
À partir de chaque réponse, formule la question en \eng{How + adjectif}, comme dans le modèle.\par
\emph{Modèle :} \eng{— It is 50 km away.} $\rightarrow$ \eng{How far is it?}
\begin{enumerate}
\item — The river is two metres deep.
\item — My grandfather is eighty years old.
\item — The wall is three metres high.
\item — The journey lasts six hours.
\item — The baby weighs four kilograms.
\item — The main road is twenty metres wide.
\end{enumerate}
\end{exercice}

\courssub{Je me prépare au Bac}

\begin{textebox}[The cleanest village in the region]
Last month, our reporter travelled to a small village near Bafoussam that has just been named the cleanest village in the West Region. Five years ago, the situation was very different: the streets were too dirty for children to play safely, and the only stream was the most polluted in the area. “How did things change?” we asked the village chief. “It was simple,” he smiled. “The young people were old enough to organise themselves, and motivated enough to work every Saturday.” They built public dustbins, planted flowers along the main road and taught the youngest pupils how important hygiene is. Today, the village has the lowest number of malaria cases in the district, and its market is the busiest and the most attractive in the area. “How far is the nearest health centre?” we asked a trader. “Only one kilometre — and now the road is good enough for ambulances,” she replied proudly. Other villages are now coming to learn from this remarkable example.
\end{textebox}

\begin{exercice}[Compréhension et langue — type Bac]
Lis le texte ci-dessus, puis réponds aux questions en anglais, par des phrases complètes.\par
\textbf{A. Comprehension questions}
\begin{enumerate}
\item Where is the village situated?
\item What was the village like five years ago? Give two details.
\item According to the chief, who changed the situation and how?
\item What shows that the village is healthier today?
\item Find in the text: (a) one superlative, (b) one structure with \eng{too...to}, (c) one structure with \eng{enough}, (d) one question with \eng{How + adjective}. Copy them.
\end{enumerate}
\textbf{B. Language work}
\begin{enumerate}
\item Transform: \eng{The streets were too dirty for children to play safely.} (Use \eng{enough}.)
\item Put into the superlative: \eng{The stream was (polluted) in the area.}
\item Correct the mistake: \eng{The young people were enough old to organise themselves.}
\end{enumerate}
\end{exercice}

\begin{exercice}[Traduction — type Bac]
Traduis le passage suivant en anglais. Tu utiliseras obligatoirement au moins deux superlatifs, une question en \eng{How + adjectif} et une structure avec \eng{too...to} ou \eng{enough...to}.\par
« Notre ville est la plus grande de la région, mais c'est aussi la plus polluée. L'air est trop sale pour les jeunes enfants. Quelle est la distance entre l'usine et l'école ? Seulement trois cents mètres. Heureusement, les élèves sont assez nombreux pour organiser une grande journée de nettoyage samedi prochain. »
\end{exercice}

\begin{exercice}[Composition — type Bac]
\textbf{Topic:} \eng{The best place to live in your region.}\par
Write a composition of \textbf{120 to 150 words} in which you describe the best place to live in your region of Cameroon. In your composition, you must:
\begin{itemize}
\item use at least \textbf{three structures} from the lesson: one superlative, one question with \eng{How + adjective} (direct or reported), and one structure with \eng{too...to} or \eng{enough...to};
\item mention the environment, health and well-being of the people who live there;
\item give your personal opinion in the conclusion.
\end{itemize}
\emph{Indications :} tu peux parler d'un quartier de Yaoundé ou de Douala, d'un village près de Buea ou de Bamenda, de la propreté, de l'air, de l'eau, des marchés, des distances (\eng{how far is the health centre?}), etc.
\end{exercice}

\begin{methode}[Réussir la composition]
\begin{enumerate}
\item \textbf{Introduction :} nomme le lieu et annonce ton opinion (\eng{In my opinion, ... is the best place to live in my region.}).
\item \textbf{Développement :} décris l'environnement (air pur, propreté, eau), la santé (centre de santé proche : \eng{How far is the health centre? Only five minutes away.}) et le bien-être (marchés, sécurité). Place un superlatif (\eng{the cleanest}, \eng{the most peaceful}) et une structure \eng{too...to} ou \eng{enough...to} (\eng{the streets are clean enough for children to play outside}).
\item \textbf{Conclusion :} donne ton avis personnel (\eng{That is why I believe...}).
\item \textbf{Relecture :} vérifie la place de \eng{enough} (après l'adjectif, avant le nom), l'article \eng{the} devant chaque superlatif, et le nombre de mots.
\end{enumerate}
\end{methode}

\newpage

\coursec{Corrigés des exercices}

\courssub{Corrigés des exercices de la leçon 34}

\begin{corrige}[Exercice 1 — QCM : la bonne forme]
\begin{enumerate}
\item B — \eng{Mount Cameroon is the highest mountain in the country.} Superlatif d'un adjectif court : \eng{the + adjectif + -est}. \eng{The most high} n'existe pas ; \eng{the higher} est un comparatif.
\item B — \eng{How long is the Wouri River? — About 160 kilometres.} \eng{How long} interroge sur la longueur ; \eng{how far} demanderait une distance entre deux lieux ; \eng{how much} une quantité ou un prix.
\item B — \eng{This water is too dirty to drink.} Structure \eng{too + adjectif + to + base verbale} : « trop sale pour être bu » (sens négatif).
\item B — \eng{She is old enough to vote.} \eng{Enough} se place \emph{après} l'adjectif : « assez âgée pour voter ».
\item B — \eng{Bamenda is colder than Yaoundé in December.} Comparatif + \eng{than} ; jamais \eng{more cold}. (Bamenda, plus haute en altitude, est plus fraîche que Yaoundé.)
\item B — \eng{How old is your grandmother? — She is seventy-five.} \eng{How old} pour l'âge.
\item A — \eng{That was the most interesting lesson of the year.} Adjectif long (trois syllabes) : \eng{the most + adjectif}, avec l'article \eng{the}.
\item B — \eng{The box is too heavy to carry.} \eng{Too + adjectif + to} ; pas de préposition \eng{for} sans complément de personne.
\item A — \eng{How tall is this mango tree? — About ten metres.} \eng{How tall} pour la hauteur d'un arbre, d'une personne ou d'un bâtiment.
\item B — \eng{He speaks English well enough to work as a guide.} \eng{Enough} se place aussi après un adverbe (\eng{well enough}).
\end{enumerate}
\textbf{Barème indicatif :} 1 point par bonne réponse correctement recopiée (10 points).
\end{corrige}

\begin{corrige}[Exercice 2 — Vrai ou faux ?]
\begin{enumerate}
\item \eng{True.} Pour demander l'âge, l'anglais utilise \eng{How old + be + sujet ?}. La structure française « combien d'années as-tu ? » ne se traduit pas mot à mot.
\item \eng{False.} \eng{Big} est un adjectif court (une syllabe) : son superlatif est \eng{the biggest}, avec redoublement du \eng{g} (consonne-voyelle-consonne). \eng{The most big} est une erreur typique de francophone.
\item \eng{True.} \eng{Too} signifie « plus qu'il n'est bon, souhaitable ou possible » : \eng{too dirty to drink} = « trop sale pour être bu ». Il ne faut pas le confondre avec \eng{very}, qui est neutre.
\item \eng{False.} \eng{Enough} se place \emph{après} l'adjectif : \eng{old enough to drive}. \eng{Enough old} est un calque du français « assez âgé ».
\item \eng{True.} \eng{How far is Buea from Douala?} interroge sur la distance entre deux lieux.
\item \eng{False.} Les adjectifs courts forment le superlatif avec \eng{-est} : \eng{the hottest} (redoublement du \eng{t}). \eng{The most} est réservé aux adjectifs longs.
\end{enumerate}
\textbf{Barème indicatif :} 1 point par réponse justifiée (6 points).
\end{corrige}

\begin{corrige}[Exercice 3 — Vocabulaire : les questions de mesure]
1 -- b ; 2 -- a ; 3 -- d ; 4 -- e ; 5 -- c ; 6 -- f.\par
Paires recopiées :
\begin{itemize}
\item \eng{How old is he?} $\rightarrow$ the age of a person.
\item \eng{How far is Buea from Douala?} $\rightarrow$ the distance between two places.
\item \eng{How tall is the building?} $\rightarrow$ the height of a thing or person.
\item \eng{How long is the bridge?} $\rightarrow$ the length of a thing.
\item \eng{How much is this fish?} $\rightarrow$ the price or quantity.
\item \eng{How deep is the lake?} $\rightarrow$ the depth of water.
\end{itemize}
Rappel : \eng{how old} (âge), \eng{how far} (distance), \eng{how tall} (hauteur), \eng{how long} (longueur ou durée), \eng{how much} (prix ou quantité indénombrable), \eng{how deep} (profondeur).
\textbf{Barème indicatif :} 1 point par paire correcte (6 points).
\end{corrige}

\begin{corrige}[Exercice 4 — Superlatifs à compléter]
\begin{enumerate}
\item \eng{The Wouri estuary is one of the largest in Central Africa.} Adjectif terminé par \eng{-e} : on ajoute seulement \eng{-st}.
\item \eng{January is often the driest month in the North of Cameroon.} Le \eng{y} devient \eng{i} devant \eng{-est}.
\item \eng{This is the best hospital in the region.} Superlatif irrégulier de \eng{good} : \eng{good -- better -- the best}.
\item \eng{The Dja Reserve is one of the oldest forests in Africa.} Adjectif court : \eng{-est}.
\item \eng{That was the worst flood the town has ever seen.} Superlatif irrégulier de \eng{bad} : \eng{bad -- worse -- the worst}.
\item \eng{My grandmother is the wisest person in our family.} \eng{Wise} se termine par \eng{-e} : \eng{-st}.
\end{enumerate}
\textbf{Point de vigilance :} dans la structure \eng{one of the + superlatif}, le nom qui suit est au \emph{pluriel} : \eng{one of the oldest forests}, « l'une des plus vieilles forêts ». Ne jamais oublier l'article \eng{the}.
\textbf{Barème indicatif :} 1 point par superlatif correctement formé avec \eng{the} (6 points).
\end{corrige}

\begin{corrige}[Exercice 5 — Traduction]
\begin{enumerate}
\item \eng{How deep is this well? — It is twelve metres deep.} (« profondeur » : nom \eng{depth}, adjectif \eng{deep}, question \eng{how deep}.)
\item \eng{Douala is the hottest town in the country in March.} (Redoublement du \eng{t} : \eng{hot -- the hottest}.)
\item \eng{This bag is too heavy for me to carry.} (« Trop lourd pour que je le porte » se rend par \eng{too heavy for me to carry} ; jamais de proposition subordonnée avec \eng{that}.)
\item \eng{She is tall enough to reach the shelf.} (\eng{Enough} après l'adjectif.)
\item \eng{How old is your little brother? — He is ten (years old).} (« Il a dix ans » se dit \eng{he is ten}, avec \eng{be}, jamais \eng{he has ten years}.)
\item \eng{It is the liveliest market in Yaoundé.} (\eng{Lively -- the liveliest} ; on peut aussi dire \eng{the busiest market}.)
\end{enumerate}
\textbf{Barème indicatif :} 2 points par phrase (structure correcte 1 point, vocabulaire et orthographe 1 point) — 12 points.
\end{corrige}

\begin{corrige}[Exercice 6 — Transformation avec \eng{enough}]
\begin{enumerate}
\item \eng{The boy is not old enough to go to school alone.}
\item \eng{This soup is not cool enough to eat now.}
\item \eng{He is not tall enough to reach the window.}
\item \eng{The road is not wide enough for two lorries to pass.}
\item \eng{She is not rested enough to walk to the market.} (accepté aussi : \eng{She is not strong enough to walk to the market.})
\end{enumerate}
\textbf{Méthode appliquée :} remplacer l'adjectif par son contraire (\eng{young} $\rightarrow$ \eng{old} ; \eng{hot} $\rightarrow$ \eng{cool} ; \eng{short} $\rightarrow$ \eng{tall} ; \eng{narrow} $\rightarrow$ \eng{wide} ; \eng{tired} $\rightarrow$ \eng{rested}), passer à la forme négative, et placer \eng{enough} après l'adjectif. Le sens reste identique : « trop jeune pour » = « pas assez âgé pour ».
\textbf{Barème indicatif :} 2 points par transformation (contraire correct 1 point, place de \eng{enough} et négation 1 point) — 10 points.
\end{corrige}

\begin{corrige}[Exercice 7 — Texte à trous : la pollution à Douala]
\begin{enumerate}
\item \eng{the most polluted} — superlatif d'un adjectif long, dans la structure \eng{one of the + superlatif}.
\item \eng{too} — \eng{the air is often too dirty for children with asthma to play outside} (« trop sale pour que les enfants... jouent »).
\item \eng{too} — \eng{the smoke is too thick for drivers to see clearly}.
\item \eng{How far} — distance entre l'usine et l'école ; majuscule en début de citation.
\item \eng{the highest} — \eng{among the highest ever recorded} : « parmi les plus élevés jamais enregistrés ».
\item \eng{the worst} — superlatif irrégulier de \eng{bad} : « la pire année ».
\item \eng{enough} — devant un \emph{nom}, \eng{enough} se place avant : \eng{enough money}.
\item \eng{How long} — durée : « combien de temps cette situation durera-t-elle ? ».
\item \eng{How much} — prix : \eng{how much does it cost?} = « combien cela coûte-t-il ? ».
\item \eng{old enough} — \eng{enough} après l'adjectif \eng{old} : « assez grande pour comprendre ».
\end{enumerate}
\textbf{Point de vigilance :} la place de \eng{enough} change selon la nature du mot qui l'accompagne : \emph{après} un adjectif ou un adverbe (\eng{old enough}, \eng{well enough}), \emph{avant} un nom (\eng{enough money}, \eng{enough time}).
\textbf{Barème indicatif :} 1 point par trou correctement complété (10 points).
\end{corrige}

\begin{corrige}[Exercice 8 — Construis la question]
\begin{enumerate}
\item \eng{How deep is the river?} (profondeur : \eng{deep}).
\item \eng{How old is your grandfather?} (âge : \eng{old}).
\item \eng{How high is the wall?} (pour un mur, une montagne, un mur d'enceinte, on préfère \eng{high} ; \eng{tall} convient aux personnes, arbres et bâtiments).
\item \eng{How long does the journey last?} (durée : \eng{long} ; auxiliaire \eng{does} avec le verbe \eng{last}).
\item \eng{How heavy is the baby?} (poids d'une personne ; accepté aussi : \eng{How much does the baby weigh?}).
\item \eng{How wide is the main road?} (largeur : \eng{wide}).
\end{enumerate}
\textbf{Point de vigilance :} avec l'auxiliaire \eng{do/does}, le verbe revient à la base verbale : \eng{does the journey last}, jamais \eng{does the journey lasts}.
\textbf{Barème indicatif :} 1 point par question correcte (6 points).
\end{corrige}

\begin{corrige}[Exercice 9 — Compréhension et langue — type Bac]
\textbf{A. Comprehension}
\begin{enumerate}
\item \eng{The village is situated near Bafoussam, in the West Region of Cameroon.}
\item \eng{Five years ago, the streets were too dirty for children to play safely, and the only stream was the most polluted in the area.}
\item \eng{According to the chief, the young people changed the situation: they organised themselves and worked every Saturday, building public dustbins, planting flowers and teaching hygiene to the youngest pupils.}
\item \eng{Today the village has the lowest number of malaria cases in the district, which shows that it is healthier.}
\item Exemples de réponses attendues : (a) \eng{the cleanest village} (ou \eng{the most polluted}, \eng{the lowest}, \eng{the busiest and the most attractive}) ; (b) \eng{the streets were too dirty for children to play safely} ; (c) \eng{old enough to organise themselves} (ou \eng{motivated enough to work}, \eng{good enough for ambulances}) ; (d) \eng{How far is the nearest health centre?}
\end{enumerate}
\textbf{B. Language work}
\begin{enumerate}
\item \eng{The streets were not clean enough for children to play safely.} (contraire de \eng{dirty} : \eng{clean} ; négation + \eng{enough} après l'adjectif).
\item \eng{The stream was the most polluted in the area.} (adjectif long : \eng{the most + participe passé}).
\item \eng{The young people were old enough to organise themselves.} (erreur corrigée : \eng{enough} se place après l'adjectif, jamais \eng{enough old}).
\end{enumerate}
\textbf{Barème indicatif :} A : 2 points par question (1 à 4), 4 points pour la question 5 (1 point par structure correctement citée) ; B : 2 points par item. Total : 20 points.
\textbf{Points de vigilance :} répondre par des phrases complètes avec verbe conjugué ; citer le texte sans le modifier ; vérifier l'article \eng{the} devant chaque superlatif.
\end{corrige}

\begin{corrige}[Exercice 10 — Traduction — type Bac]
\textbf{Proposition de traduction :}\par
\eng{Our town is the largest in the region, but it is also the most polluted. The air is too dirty for young children. How far is the factory from the school? Only three hundred metres. Fortunately, there are enough pupils to organise a big clean-up day next Saturday.}\par
\textbf{Justifications :}
\begin{itemize}
\item « la plus grande » : \eng{the largest} (adjectif terminé par \eng{-e} : \eng{-st}) ; « la plus polluée » : \eng{the most polluted} (adjectif long). Les deux superlatifs imposés sont présents, chacun avec \eng{the}.
\item « trop sale pour » : \eng{too dirty for young children} — structure \eng{too + adjectif + for + nom}.
\item « Quelle est la distance entre l'usine et l'école ? » : \eng{How far is the factory from the school?} — jamais \eng{What is the distance...} dans une question naturelle de ce niveau.
\item « assez nombreux pour organiser » : \eng{there are enough pupils to organise} — \eng{enough} devant le nom pluriel \eng{pupils}.
\item « trois cents mètres » : \eng{three hundred metres} — \eng{hundred} reste invariable après un nombre.
\end{itemize}
\textbf{Barème indicatif :} 10 points — respect des quatre structures imposées (4 points), exactitude grammaticale (3 points), vocabulaire et orthographe (2 points), fluidité (1 point).
\textbf{Points de vigilance :} pas de \eng{s} à \eng{hundred} ; \eng{enough} avant le nom mais après l'adjectif ; article \eng{the} devant chaque superlatif.
\end{corrige}

\begin{corrige}[Exercice 11 — Composition — type Bac]
\textbf{Proposition de rédaction modèle (environ 135 mots) :}\par
\eng{In my opinion, Molyko is the best place to live in the South-West Region. It is one of the cleanest neighbourhoods of Buea, and the air there is much fresher than in the big industrial cities.}\par
\eng{The people of Molyko enjoy a healthy environment. The streets are clean enough for children to play outside safely, and the water from the springs is never too dirty to use. How far is the nearest health centre? Only ten minutes on foot, which is very important for sick people and old people. The market is the busiest in the town, so families can always find fresh food at good prices.}\par
\eng{Life is also pleasant because the climate is neither too hot nor too cold. The inhabitants are friendly and motivated enough to keep their quarter beautiful.}\par
\eng{That is why I believe Molyko is the most comfortable place to live in my region.}\par
\textbf{Barème indicatif :} 20 points — respect des consignes et des trois structures imposées (6 points), correction grammaticale (6 points), richesse du vocabulaire lié à l'environnement et à la santé (4 points), organisation et cohérence (2 points), respect du nombre de mots (2 points).\par
\textbf{Points de vigilance :}
\begin{itemize}
\item Vérifier que chaque superlatif est précédé de \eng{the} : \eng{the cleanest}, \eng{the busiest}, \eng{the most comfortable}.
\item Placer \eng{enough} après l'adjectif (\eng{clean enough}) et avant le nom (\eng{enough time}).
\item Ne pas traduire « il y a » par \eng{it has} : utiliser \eng{there is / there are}.
\item Éviter le calque « les gens sont assez motivés » $\rightarrow$ \eng{motivated enough}, jamais \eng{enough motivated}.
\end{itemize}
\end{corrige}

\newpage

\coursec{Fiche bilan}

\begin{popfigure}
\begin{tikzpicture}[mindmap, grow cyclic, every node/.style={concept, font=\small}, root concept/.append style={concept color=popblue, font=\bfseries\small, minimum size=3.2cm, text=white}, level 1 concept/.append style={level distance=3.6cm, sibling angle=60, font=\footnotesize, minimum size=2.3cm}]
\node[root concept, align=center]{Adjectives:\\difficult\\structures}
  child[concept color=poporange, text=black]{ node[align=center]{How + adj\\questions de\\mesure} }
  child[concept color=popgreen, text=black]{ node[align=center]{Superlatif\\the + -est\\the most...} }
  child[concept color=popgold, text=black]{ node[align=center]{too + adj\\(+ to + BV)\\= trop} }
  child[concept color=poppurple, text=white]{ node[align=center]{adj + enough\\(+ to + BV)\\= assez} }
  child[concept color=poppink, text=black]{ node[align=center]{adj + to + BV\\easy, difficult\\glad, sorry} }
  child[concept color=popmuted, text=white]{ node[align=center]{Pièges\\francophones\\ordre des mots} };
\end{tikzpicture}
\legende{Carte mentale de la leçon : les structures difficiles avec les adjectifs}
\end{popfigure}

\begin{bilan}
\textbf{1. Lexique clé (English / Français)}
\begin{itemize}
  \item \eng{deep / depth} : profond / la profondeur ; \eng{wide / width} : large / la largeur ; \eng{high / height} : haut / la hauteur ; \eng{long / length} : long / la longueur.
  \item \eng{healthy / unhealthy} : sain / malsain ; \eng{polluted} : pollué ; \eng{harmful} : nocif ; \eng{safe} : sûr, sans danger.
  \item \eng{enough} : assez ; \eng{too} : trop ; \eng{the most / the least} : le plus / le moins.
\end{itemize}

\textbf{2. Règles à connaître par cœur}
\begin{center}
\begin{tabularx}{\linewidth}{|X|X|X|}
\hline
\rowcolor{popblueL} Structure & Emploi & Exemple \\
\hline
\eng{How + adj + is/are...?} & poser une question de mesure & \eng{How deep is the river?} « Quelle est la profondeur de la rivière ? » \\
\hline
\eng{the + adj + -est} (1 syllabe) & superlatif de supériorité & \eng{Mount Cameroon is the highest peak in West Africa.} « Le mont Cameroun est le sommet le plus haut d'Afrique de l'Ouest. » \\
\hline
\eng{the most + adj} (longs) & superlatif des adjectifs longs & \eng{This is the most polluted area of Douala.} « C'est la zone la plus polluée de Douala. » \\
\hline
\eng{too + adj (+ to + BV)} & excès, limite négative & \eng{The water is too dirty to drink.} « L'eau est trop sale pour être bue. » \\
\hline
\eng{adj + enough (+ to + BV)} & suffisance & \eng{She is old enough to understand.} « Elle est assez grande pour comprendre. » \\
\hline
\eng{adj + to + BV} & réaction, facilité & \eng{I am glad to help.} « Je suis heureux d'aider. » \\
\hline
\end{tabularx}
\end{center}

\textbf{3. Points de forme irréguliers}
\begin{itemize}
  \item Superlatifs irréguliers : \eng{good} $\to$ \eng{the best} ; \eng{bad} $\to$ \eng{the worst} ; \eng{far} $\to$ \eng{the farthest/furthest}.
  \item Adjectifs en \eng{-y} : \eng{healthy} $\to$ \eng{the healthiest} ; consonne double : \eng{big} $\to$ \eng{the biggest}, \eng{hot} $\to$ \eng{the hottest}.
  \item Après le superlatif : \eng{in} + lieu (\eng{the best student in the class}), jamais \eng{of} pour un lieu.
\end{itemize}

\textbf{4. Méthode express}
\begin{enumerate}
  \item Question de mesure ? $\to$ \eng{How + adjectif} (jamais \eng{What}).
  \item Comparer un élément à tout un groupe ? $\to$ superlatif avec \eng{the}, adjectif court = \eng{-est}, adjectif long = \eng{the most}.
  \item Idée de « trop » $\to$ \eng{too} AVANT l'adjectif ; idée de « assez » $\to$ \eng{enough} APRÈS l'adjectif.
  \item « pour + infinitif » consécutif $\to$ \eng{to + base verbale} après \eng{too...} ou \eng{...enough}.
\end{enumerate}
\end{bilan}

\begin{aretenir}[Checklist avant le Bac]
\begin{itemize}
  \item[$\square$] Je sais poser une question de mesure avec \eng{How + adjective} : \eng{How high is Mount Cameroon?}
  \item[$\square$] Je place toujours \eng{the} devant le superlatif : \eng{the cleanest}, \eng{the most dangerous}.
  \item[$\square$] Je connais les superlatifs irréguliers : \eng{the best}, \eng{the worst}, \eng{the farthest}.
  \item[$\square$] Je double la consonne finale quand il le faut : \eng{the biggest}, \eng{the hottest}, \eng{the thinnest}.
  \item[$\square$] Je transforme le \eng{-y} en \eng{-iest} : \eng{the healthiest}, \eng{the happiest}.
  \item[$\square$] Je mets \eng{too} AVANT l'adjectif : \eng{too expensive}, jamais \eng{expensive too}.
  \item[$\square$] Je mets \eng{enough} APRÈS l'adjectif mais AVANT le nom : \eng{strong enough}, \eng{enough water}.
  \item[$\square$] Je complète avec \eng{to + base verbale} : \eng{too weak to walk}, \eng{old enough to vote}.
  \item[$\square$] J'utilise \eng{in} (et non \eng{of}) après un superlatif suivi d'un lieu : \eng{the tallest boy in the school}.
  \item[$\square$] Je ne traduis jamais « Quelle est la profondeur... ? » par \eng{What is the deep...?} mais par \eng{How deep is...?}
  \item[$\square$] Je sais transformer une phrase affirmative en structure avec \eng{too...to} ou \eng{enough...to}.
  \item[$\square$] Je vérifie l'accord du verbe après \eng{How + adj} : \eng{How tall are you?} / \eng{How tall is he?}
\end{itemize}
\end{aretenir}

\findecours

\end{document}
